% Plantilla del curso intensivo de Lean. NO editar CursoLean.tex: se genera con
%   python3 reports/curso-lean/build.py
% Las directivas %%CODE, %%EXCERPT, %%STATE, %%STEP y %%LEAN se sustituyen por
% código real del repositorio y por estados de prueba exportados por Lean.
\documentclass[11pt,a4paper]{article}
\usepackage[a4paper,margin=2.3cm,headheight=14pt]{geometry}
\usepackage{fontspec}
\usepackage{unicode-math}
\usepackage[spanish,es-noshorthands,es-tabla]{babel}
\directlua{
  luaotfload.add_fallback("leanfallback", {
    "STIX Two Math:mode=harf;",
    "DejaVu Sans:mode=harf;",
    "Apple Symbols:mode=harf;",
  })
}
\setmainfont{TeX Gyre Pagella}
\setmathfont{TeX Gyre Pagella Math}
\setsansfont{DejaVu Sans}[Scale=MatchLowercase]
\setmonofont{DejaVu Sans Mono}[Scale=MatchLowercase, RawFeature={fallback=leanfallback}]
\usepackage{xcolor}
\usepackage{fvextra}
\usepackage[most]{tcolorbox}
\usepackage{booktabs,tabularx,array,longtable}
\usepackage{bussproofs}
\usepackage{enumitem}
\usepackage{fancyhdr}
\usepackage[hidelinks]{hyperref}
\usepackage{microtype}

\definecolor{kw}{HTML}{1F4E9C}
\definecolor{tac}{HTML}{8A2B8F}
\definecolor{com}{HTML}{5E6B5E}
\definecolor{fondocodigo}{HTML}{F7F8FA}
\definecolor{bordecodigo}{HTML}{B8C2CF}
\definecolor{fondoestado}{HTML}{FFF9EC}
\definecolor{bordeestado}{HTML}{D9B45A}
\definecolor{fondoesquema}{HTML}{EEF5EE}
\definecolor{bordeesquema}{HTML}{5E9A62}
\definecolor{fondonota}{HTML}{F1EEF8}
\definecolor{bordenota}{HTML}{8C7DB8}

\newcommand{\LKw}[1]{\textcolor{kw}{\bfseries #1}}
\newcommand{\LTac}[1]{\textcolor{tac}{\bfseries #1}}
\newcommand{\LCom}[1]{\textcolor{com}{\itshape #1}}
\newcommand{\LCont}[1]{\textcolor{gray}{#1}}
\directlua{
  function leanbreak(s)
    local previo = ""
    for _, c in utf8.codes(s) do
      local ch = utf8.char(c)
      if ch:match("^[A-Z]$") and previo:match("^[a-z]$") then
        tex.sprint("\\allowbreak{}")
      end
      tex.sprint(-2, ch)
      if ch == "_" or ch == "." then tex.sprint("\\allowbreak{}") end
      previo = ch
    end
  end
}
\newcommand{\lean}[1]{\texttt{\directlua{leanbreak("\luaescapestring{\detokenize{#1}}")}}}
\newcommand{\ruta}[2]{\texttt{#1}\,:\,#2}

\fvset{fontsize=\footnotesize, baselinestretch=1.05}

\newtcolorbox{codigolean}[1]{enhanced, breakable, colback=fondocodigo,
  colframe=bordecodigo, boxrule=0.5pt, arc=2pt, left=5pt, right=5pt,
  top=3pt, bottom=3pt, colbacktitle=bordecodigo!35, coltitle=black,
  fonttitle=\sffamily\scriptsize, title={#1}, toptitle=1pt, bottomtitle=1pt,
  before skip=6pt, after skip=8pt}
\newtcolorbox{estadolean}[1]{enhanced, breakable, colback=fondoestado,
  colframe=bordeestado, boxrule=0.5pt, arc=2pt, left=5pt, right=5pt,
  top=3pt, bottom=3pt, colbacktitle=bordeestado!30, coltitle=black,
  fonttitle=\sffamily\scriptsize, title={#1}, toptitle=1pt, bottomtitle=1pt,
  before skip=4pt, after skip=4pt}
\newtcolorbox{esquema}[1]{enhanced, breakable, colback=fondoesquema,
  colframe=bordeesquema, boxrule=0.6pt, arc=3pt,
  colbacktitle=bordeesquema, coltitle=white, fonttitle=\sffamily\bfseries,
  title={#1}, before skip=10pt, after skip=10pt, before upper={\raggedright}}
\newtcolorbox{nota}[1][Nota]{enhanced, breakable, colback=fondonota,
  colframe=bordenota, boxrule=0.5pt, arc=3pt, colbacktitle=bordenota,
  coltitle=white, fonttitle=\sffamily\bfseries\small, title={#1},
  before skip=8pt, after skip=8pt}
\newcommand{\tacticaentre}[1]{\begin{center}\vspace{-2pt}\(\Downarrow\)\enspace
  \texttt{\small #1}\vspace{-2pt}\end{center}}

\newcommand{\forma}{\par\noindent\textbf{Forma matemática.}\enspace}
\newcommand{\lectura}{\par\smallskip\noindent\textbf{Cómo se lee.}\enspace}
\newcommand{\donde}{\par\smallskip\noindent\textbf{Dónde aparece.}\enspace}

\pagestyle{fancy}
\fancyhf{}
\fancyhead[L]{\small\sffamily Curso intensivo de Lean~4 a través de la tesis}
\fancyhead[R]{\small\sffamily\nouppercase{\leftmark}}
\fancyfoot[C]{\small\thepage}
\renewcommand{\headrulewidth}{0.4pt}

\setlist{itemsep=2pt, topsep=4pt}
\setlength{\parskip}{4pt}
\setlength{\parindent}{0pt}
\EnableBpAbbreviations

\title{\textbf{Curso intensivo de Lean~4 a través de la tesis}\\[4pt]
  \large Esquemas de demostración: cómo se interpretan y cómo se prueban}
\author{Material de apoyo para explicar las formalizaciones de\\
  \texttt{Thesis/Prop}, \texttt{Thesis/Sort} y \texttt{Thesis/Teoria}}
\date{}

\begin{document}
\maketitle
\thispagestyle{empty}

\begin{abstract}
\noindent
Este documento es una guía de lectura del código Lean de la tesis. No es un
manual general de Lean: solo se explican las construcciones que aparecen en
las demostraciones de los capítulos 7 (cálculo~\(\lambda\) y tipos simples),
8 (cubo~\(\lambda\)), 9 (deducción natural), 11 (completitud) y 12
(Quicksort). El núcleo del documento es un catálogo de \emph{esquemas de
demostración}. Cada esquema enlaza una forma de razonamiento matemático con la
plantilla Lean que la implementa, la lectura en español de esa plantilla y los
lugares de la tesis donde se usa.

Los fragmentos de código no están copiados a mano. Un script los extrae del
repositorio (se indica archivo y línea). Los recuadros amarillos muestran
\emph{estados de prueba reales}, exportados por Lean~4.29 al compilar el
proyecto (\texttt{artifacts/explorer/*.json}). En esos recuadros se omiten las
hipótesis inaccesibles que Lean marca con~\texttt{✝}.
\end{abstract}

\tableofcontents
\newpage

%=====================================================================
\section{Cómo usar este documento}
%=====================================================================

La idea central es esta: \textbf{una demostración en Lean es la misma
demostración que en papel, escrita con un vocabulario restringido y
comprobada por una máquina}. Cada paso del papel («sea \(x\)», «supongamos
\(P\)», «por hipótesis de inducción», «distinguimos dos casos») corresponde a
una táctica o a un término de Lean. Si se conoce esa correspondencia, la prueba
formal se lee como un texto matemático, y la prueba en papel se puede
reconstruir a partir del código.

Para preparar la explicación de un teorema concreto se sugiere este orden:
\begin{enumerate}
  \item Leer el \textbf{enunciado} y traducirlo a lenguaje matemático (sección~\ref{sec:basicos}).
  \item Leer los comentarios \texttt{A probar} y \texttt{Método}, que en esta tesis encabezan casi todas las pruebas.
  \item Identificar el \textbf{esquema} principal: inducción, casos, contradicción, cadena de igualdades, construcción de una derivación\ldots{} (sección~\ref{sec:esquemas}).
  \item Recorrer cada caso o rama con la tabla de tácticas (sección~\ref{sec:tacticas}).
  \item Separar los pasos de contenido matemático de los pasos de «contabilidad» (pertenencia a un contexto, reescrituras de definiciones, condiciones laterales).
\end{enumerate}

La sección~\ref{sec:mapa} da un mapa de toda la tesis: para cada resultado
formalizado indica el archivo, el teorema y los esquemas que usa. La
sección~\ref{sec:explicar} contiene una receta para exponer una prueba en voz
alta y una lista de preguntas frecuentes.

%=====================================================================
\section{Lo básico: enunciados, pruebas y estados}\label{sec:basicos}
%=====================================================================

\subsection{Proposiciones como tipos}

Lean es un lenguaje de programación con tipos dependientes y, a la vez, un
verificador de demostraciones. Se basa en el \emph{cálculo de construcciones
inductivas}, una extensión del sistema \(\lambda C\) que la tesis presenta como
vértice del cubo~\(\lambda\) (capítulo~8). Por eso la correspondencia de
Curry--Howard de los capítulos 7 y~8 no es aquí una analogía, sino el
funcionamiento real del sistema:

\begin{center}
\begin{tabular}{@{}ll@{}}
\toprule
\textbf{Matemática} & \textbf{Lean} \\
\midrule
proposición \(P\) & tipo \lean{P : Prop} \\
demostración de \(P\) & término \lean{h : P} \\
\(P \Rightarrow Q\) & tipo de funciones \lean{P → Q} \\
\(\forall x \in A,\ Q(x)\) & tipo de funciones dependientes \lean{∀ x : A, Q x} \\
aplicar una implicación (modus ponens) & aplicar una función: \lean{h hp} \\
comprobar una prueba & comprobar que un término tiene el tipo pedido \\
\bottomrule
\end{tabular}
\end{center}

Un teorema es una definición cuyo tipo es una proposición:
\begin{codigolean}{Anatomía de un teorema (esquema general)}
\begin{Verbatim}[commandchars=\\\{\}]
\LKw{theorem} nombre (hipótesis₁ : A) \{implícita : B\} : enunciado := \LKw{by}
  \LCom{-- A probar: …  (comentario de la tesis)}
  \LCom{-- Método: …}
  táctica₁
  táctica₂
\end{Verbatim}
\end{codigolean}

Los argumentos entre paréntesis \lean{( )} son explícitos: quien usa el teorema
los proporciona. Los argumentos entre llaves \lean{{ }} son implícitos: Lean
los deduce del contexto. Lo que sigue a los dos puntos es el enunciado; lo que
sigue a \lean{:=} es la demostración. La palabra \lean{by} abre el
\emph{modo táctico}: en lugar de escribir el término de prueba completo, se dan
instrucciones que lo construyen paso a paso.

\subsection{El estado de prueba}

En modo táctico Lean mantiene un \emph{estado}: una lista de \emph{metas}
(lo que falta probar). Cada meta tiene un contexto de hipótesis disponibles,
escritas arriba como \lean{nombre : tipo}, y una conclusión, escrita después
del símbolo~\lean{⊢}. Por ejemplo, al empezar la prueba de la introducción de
la doble negación (capítulo~11) el estado es:

\begin{estadolean}{Estado antes de \texttt{let Δ := insert (\textasciitilde{}φ) Γ}}
\begin{Verbatim}[commandchars=\\\{\}]
Γ : Set Formula
φ : Formula
d : ND Γ φ
⊢ ND Γ (~~φ)
\end{Verbatim}
\end{estadolean}

Se lee así: «tenemos un contexto \(\Gamma\), una fórmula \(\varphi\) y una
derivación \(d\) de \(\Gamma \vdash \varphi\); hay que construir una derivación
de \(\Gamma \vdash \neg\neg\varphi\)». Cada táctica transforma este estado.
Algunas añaden hipótesis (\lean{have}, \lean{intro}), otras cambian la meta
(\lean{rw}, \lean{change}), otras la dividen en varias (\lean{constructor},
\lean{rcases}, \lean{induction}) y otras la cierran (\lean{exact}). La prueba
termina cuando no quedan metas.

\begin{nota}[Idea clave para la exposición]
Al explicar una prueba Lean conviene describir el \emph{estado}, no la táctica.
«Después de \lean{intro h} tenemos \(h : P\) y debemos probar \(Q\)» se
entiende mejor que «aplicamos intro». Los recuadros amarillos de este documento
muestran esos estados tal como Lean los calcula.
\end{nota}

\subsection{Modo término y modo táctico}

Las dos formas producen el mismo objeto: un término cuyo tipo es el
enunciado. El modo término es conveniente cuando la prueba es una sola
aplicación de lemas ya demostrados:

\begin{codigolean}{\ruta{Thesis/Prop/Main.lean}{20}}
\begin{Verbatim}[commandchars=\\\{\}]
\LKw{theorem} identidad_por_completitud : ND (∅ : Set Formula) tautologiaIdentidad :=
  completitud_ND tautologiaIdentidad identidad_es_tautologia
\end{Verbatim}
\end{codigolean}
\begin{codigolean}{\ruta{Thesis/Teoria/DeduccionNatural.lean}{253}}
\begin{Verbatim}[commandchars=\\\{\}]
\LKw{theorem} contrapositiva (p q : Formula) : (p →ₚ q) ⊣⊢ₚ (¬ₚq →ₚ ¬ₚp) :=
  ⟨contrapositiva_ida p q, contrapositiva_vuelta p q⟩
\end{Verbatim}
\end{codigolean}

En el segundo ejemplo, \lean{⟨a, b⟩} es el \emph{constructor anónimo}: produce
una prueba de una conjunción a partir de las pruebas de sus dos partes
(aquí \lean{⊣⊢ₚ} está definido como una conjunción de dos secuentes).

\subsection{Tipos inductivos: datos y reglas}\label{sec:inductivos}

Casi todo lo que la tesis formaliza se define con \lean{inductive}. Hay dos
usos distintos de esta palabra, y conviene distinguirlos al explicar.

\textbf{Tipos de datos inductivos.} Son la sintaxis: fórmulas, términos,
tipos, listas. Cada constructor es una forma de construir un elemento.
\begin{codigolean}{\ruta{Thesis/Prop/Syntax.lean}{3}}
\begin{Verbatim}[commandchars=\\\{\}]
\LKw{inductive} Formula : Type
| atom : String → Formula
| neg  : Formula → Formula
| impl : Formula → Formula → Formula
\end{Verbatim}
\end{codigolean}

\textbf{Predicados inductivos.} Son relaciones definidas por reglas:
derivabilidad, juicios de tipado, pasos de reducción. Cada constructor es una
\emph{regla de inferencia}: sus argumentos son las premisas y su tipo de
retorno es la conclusión. La relación es la \emph{menor} cerrada bajo esas
reglas, y por eso admite inducción sobre derivaciones (esquema~11, sección~\ref{esq:induccion-derivaciones}).
\begin{codigolean}{\ruta{Thesis/Prop/NaturalDeduction.lean}{37}}
\begin{Verbatim}[commandchars=\\\{\}]
\LKw{inductive} ND : Set Formula → Formula → Prop
| hyp \{Γ φ\} : φ ∈ Γ → ND Γ φ
| impI \{Γ φ ψ\} : ND (insert φ Γ) ψ → ND Γ (φ ⟶ ψ)
| impE \{Γ φ ψ\} : ND Γ (φ ⟶ ψ) → ND Γ φ → ND Γ ψ
| negI \{Γ φ\} : ND (insert φ Γ) Falsedad → ND Γ (~φ)
| negE \{Γ φ\} : ND Γ (~φ) → ND Γ φ → ND Γ Falsedad
| botE \{Γ φ\} : ND Γ Falsedad → ND Γ φ
\LCom{/-- Por RAA: de Γ, ¬φ ⊢ ⊥ se concluye Γ ⊢ φ, descargando ¬φ. -/}
| classical \{Γ φ\} : ND (insert (~φ) Γ) Falsedad → ND Γ φ
\end{Verbatim}
\end{codigolean}

Comparando con la presentación en papel:
\begin{center}
\renewcommand{\arraystretch}{2.4}
\begin{tabular}{@{}cc@{}}
\AXC{$\varphi\in\Gamma$}\RightLabel{\scriptsize\texttt{hyp}}\UIC{$\Gamma\vdash\varphi$}\DisplayProof
&
\AXC{$\Gamma,\varphi\vdash\psi$}\RightLabel{\scriptsize\texttt{impI}}\UIC{$\Gamma\vdash\varphi\to\psi$}\DisplayProof
\\
\AXC{$\Gamma\vdash\varphi\to\psi$}\AXC{$\Gamma\vdash\varphi$}\RightLabel{\scriptsize\texttt{impE}}\BIC{$\Gamma\vdash\psi$}\DisplayProof
&
\AXC{$\Gamma,\varphi\vdash\bot$}\RightLabel{\scriptsize\texttt{negI}}\UIC{$\Gamma\vdash\neg\varphi$}\DisplayProof
\\
\AXC{$\Gamma\vdash\neg\varphi$}\AXC{$\Gamma\vdash\varphi$}\RightLabel{\scriptsize\texttt{negE}}\BIC{$\Gamma\vdash\bot$}\DisplayProof
&
\AXC{$\Gamma\vdash\bot$}\RightLabel{\scriptsize\texttt{botE}}\UIC{$\Gamma\vdash\varphi$}\DisplayProof
\quad
\AXC{$\Gamma,\neg\varphi\vdash\bot$}\RightLabel{\scriptsize\texttt{classical}}\UIC{$\Gamma\vdash\varphi$}\DisplayProof
\end{tabular}
\end{center}
La lectura \(\Gamma, \varphi\) del papel es \lean{insert φ Γ} en Lean: el
contexto es un conjunto \lean{Set Formula}. Lo mismo ocurre con los otros
cálculos de la tesis: \lean{Derivable} (capítulo~9, con todos los conectivos),
\lean{Deriva} (tipos simples, \(\lambda2\) y \(\lambda\underline{\omega}\)),
\lean{PasoBeta}, \lean{Cadena} y \lean{AlfaTesis}.

\subsection{Nivel objeto y metanivel}\label{sec:objeto-meta}

Las formalizaciones lógicas de la tesis tienen dos lenguajes superpuestos, y
mezclarlos es la fuente de confusión más frecuente:

\begin{center}
\small
\begin{tabularx}{\textwidth}{@{}l>{\raggedright\arraybackslash}X>{\raggedright\arraybackslash}X@{}}
\toprule
& \textbf{Lógica objeto (lo que se estudia)} & \textbf{Metalógica (Lean)} \\
\midrule
negación & \lean{~φ} (\lean{Formula.neg}), \lean{¬ₚφ} & \lean{¬ P} \\
implicación & \lean{φ ⟶ ψ}, \lean{φ →ₚ ψ} & \lean{P → Q} \\
falsedad & \lean{Falsedad}, \lean{⊥ₚ} & \lean{False} \\
«es demostrable» & \lean{ND Γ φ}, \lean{Γ ⊢ₚ φ} & tener un término de ese tipo \\
reducción al absurdo & \lean{ND.classical}, \lean{Derivable.raa} & \lean{by_contra}, \lean{Classical.byContradiction} \\
análisis por casos & \lean{analisis_por_casos}, \lean{Derivable.disyE} & \lean{by_cases}, \lean{rcases} \\
\bottomrule
\end{tabularx}
\end{center}

Los comentarios de la tesis marcan esta distinción cuando importa («Por RAA
metateórica\ldots», «Introducción de doble negación objeto; no es la
eliminación anterior»). En el Lema de Kalmár, por ejemplo, se usa la
eliminación de la doble negación \emph{en la metalógica} para decidir el valor
de verdad de \(p\), y la introducción de la doble negación \emph{en el cálculo
ND} para construir la derivación.

\begin{nota}[Cuidado con \texttt{\textasciitilde}]
En \texttt{Thesis/Prop} el prefijo \lean{~} es la negación de fórmulas. En
\texttt{Thesis/Sort} el infijo \lean{l₁ ~ l₂} es \lean{List.Perm}: «\(l_1\) es
una permutación de \(l_2\)». Son notaciones distintas en archivos distintos.
\end{nota}

\subsection{Convenciones de escritura de la tesis}

\begin{itemize}
  \item Los identificadores propios están en español, por ejemplo \lean{debilitamiento},
    \lean{quicksort_ordenada} o \lean{contextoLiterales}. Las tácticas y las
    bibliotecas externas conservan sus nombres en inglés.
  \item Casi todas las pruebas empiezan con dos comentarios:
    \emph{A probar} (la meta en palabras) y \emph{Método} (el esquema que se usará).
  \item Prefijos de nombres: \lean{h…} para hipótesis proposicionales
    (\lean{hsub}, \lean{hCasos}); \lean{d…} para derivaciones del cálculo
    (\lean{dNeg}, \lean{dFalsedad}); \lean{hipInd…} para hipótesis de
    inducción; \lean{l1}, \lean{l2}, \ldots{} para las líneas de una derivación
    en cajas; \lean{paso1}, \lean{paso2}, \ldots{} para los pasos numerados de
    los ejemplos de tipado.
  \item Se prefiere el estilo \emph{hacia adelante}: una secuencia de
    \lean{have} con su tipo escrito, de modo que cada línea enuncia un hecho
    intermedio legible. Esto hace las pruebas más largas que el mínimo, pero
    cada paso se puede leer sin ejecutar Lean.
  \item Los pasos automáticos (\lean{simp}, \lean{decide}, \lean{omega}) se
    reservan para condiciones laterales decidibles o para desplegar
    definiciones; no hacen el razonamiento principal.
\end{itemize}

%=====================================================================
\section{Conectivos: introducción y eliminación}\label{sec:conectivos}
%=====================================================================

En Lean, como en deducción natural, cada conectivo tiene una forma de
\emph{probarlo} (introducción) y una forma de \emph{usarlo} (eliminación). Esta
tabla es la base de todos los esquemas siguientes.

\begin{center}
\small
\renewcommand{\arraystretch}{1.25}
\begin{tabularx}{\textwidth}{@{}l>{\raggedright\arraybackslash}X>{\raggedright\arraybackslash}X@{}}
\toprule
\textbf{Conectivo} & \textbf{Para probarlo (meta)} & \textbf{Para usarlo (hipótesis \lean{h})} \\
\midrule
\(P \to Q\) & \lean{intro hp}, luego probar \(Q\) & \lean{h hp : Q} (aplicar), o \lean{apply h} \\
\(\forall x, Q\,x\) & \lean{intro x} & \lean{h a : Q a} \\
\(P \wedge Q\) & \lean{constructor}, \lean{⟨hp, hq⟩}, \lean{And.intro}, \lean{refine ⟨?_, ?_⟩} & \lean{h.left}, \lean{h.right}, \lean{obtain ⟨hp, hq⟩ := h} \\
\(P \vee Q\) & \lean{Or.inl hp}, \lean{Or.inr hq} & \lean{rcases h with hp | hq} (dos metas) \\
\(P \leftrightarrow Q\) & \lean{constructor} y dos implicaciones & \lean{h.mp : P → Q}, \lean{h.mpr : Q → P}, \lean{rw [h]} \\
\(\neg P\) (es \(P \to \mathrm{False}\)) & \lean{intro hp}, luego probar \lean{False} & \lean{h hp : False} \\
\(\mathrm{False}\) & (no se introduce) & \lean{False.elim h}, \lean{absurd hp hnp}, \lean{exact (h hp).elim} \\
\(\exists x, Q\,x\) & \lean{⟨a, prueba⟩} (dar un testigo) & \lean{obtain ⟨x, hx⟩ := h} \\
\(a = b\) & \lean{rfl}, \lean{rw}, \lean{calc}, \lean{simp} & \lean{rw [h]}, \lean{subst}, \lean{h ▸ e} \\
\(x \in \{y \mid P\,y\}\) & probar \(P\,x\) (es la misma proposición) & usarla como \(P\,x\) \\
\bottomrule
\end{tabularx}
\end{center}

\begin{nota}[Igualdad por definición]
Lean identifica expresiones que son iguales \emph{por cómputo de las
definiciones} (igualdad definicional). Por eso
\lean{have hNoIdentidad : ¬ (v "⊥" → v "⊥") := h} es válido cuando
\lean{h : evaluar v Falsedad}: al desplegar \lean{evaluar} y \lean{Falsedad}
se obtiene exactamente ese tipo. Las tácticas \lean{change} y \lean{show}
reemplazan la meta por otra definicionalmente igual. No cambian lo que se
prueba: solo lo escriben de otra forma.
\end{nota}

%=====================================================================
\section{Vocabulario de tácticas usado en la tesis}\label{sec:tacticas}
%=====================================================================

\begin{center}
\small
\renewcommand{\arraystretch}{1.2}
\begin{longtable}{@{}>{\raggedright\arraybackslash}p{0.27\textwidth}>{\raggedright\arraybackslash}p{0.36\textwidth}>{\raggedright\arraybackslash}p{0.31\textwidth}@{}}
\toprule
\textbf{Táctica} & \textbf{Efecto en el estado} & \textbf{Frase en papel} \\
\midrule
\endhead
\lean{intro x h} & mueve \(\forall x\) y premisas de la meta al contexto & «Sea \(x\) y supongamos \(h\)» \\
\lean{exact t} & cierra la meta con el término \lean{t} & «lo que se quería, por \ldots» \\
\lean{have h : P := t} & añade el hecho \(P\) (probado por \lean{t}) & «tenemos \(P\) por \ldots» \\
\lean{have h : P := by …} & idem, con una subprueba & «afirmamos \(P\); en efecto, \ldots» \\
\lean{let x := e} & introduce una abreviatura local & «sea \(x := e\)» \\
\lean{apply f} & reduce la meta a las premisas de \lean{f} & «basta probar las hipótesis de \ldots» \\
\lean{refine ⟨?_, ?_⟩} & divide la meta en las piezas marcadas \lean{?_} & «probamos cada parte» \\
\lean{constructor} & divide \(\wedge\), \(\leftrightarrow\) en sus partes & «probamos ambas direcciones» \\
\lean{rcases h with a | b}, \lean{obtain ⟨a, b⟩ := h}, \lean{rintro} & descompone \(\vee\), \(\wedge\), \(\exists\) & «distinguimos casos» / «sea \(x\) tal que» \\
\lean{cases h with | c …} & análisis por constructores de \lean{h} & «\(h\) solo pudo obtenerse por la regla \ldots» \\
\lean{induction x with | c …} & una meta por constructor, con hipótesis inductivas & «por inducción sobre \ldots» \\
\lean{by_cases h : P} & dos metas: con \(h : P\) y con \(h : \neg P\) & «si \(P\)\ldots; si no\ldots» (tercero excluido) \\
\lean{by_contra h} & supone la negación de la meta; la meta pasa a ser \lean{False} & «supongamos lo contrario» \\
\lean{rw [e₁, e₂] (at h)} & reescribe con igualdades o equivalencias & «sustituyendo \ldots» \\
\lean{subst h} & elimina una variable igualada por \lean{h} & «como \(x = a\), \ldots» \\
\lean{change P}, \lean{show P} & reescribe la meta en forma definicionalmente igual & «es decir, hay que probar \(P\)» \\
\lean{calc a = b := … _ = c := …} & cadena de igualdades o relaciones transitivas & «\(a = b = c\)» \\
\lean{simp only [l]}, \lean{simpa … using h} & simplificación con los lemas indicados & «desplegando las definiciones» \\
\lean{unfold f} & despliega la definición \lean{f} & «por definición de \(f\)» \\
\lean{ext x} & reduce la igualdad de conjuntos a pertenencia & «por extensionalidad» \\
\lean{simp}, \lean{decide}, \lean{omega} & cierran condiciones laterales decidibles o aritméticas & «trivialmente», «por cálculo» \\
\lean{termination_by m}, \lean{decreasing_by} & declaran y prueban que la medida \lean{m} decrece & «inducción sobre la longitud» \\
\lean{· …} & enfoca la siguiente meta pendiente & «en el primer caso\ldots» \\
\bottomrule
\end{longtable}
\end{center}

Términos de uso frecuente: \lean{False.elim} (de lo falso, cualquier cosa),
\lean{absurd hp hnp}, \lean{Or.inl}/\lean{Or.inr}, \lean{And.intro},
\lean{h.mp}/\lean{h.mpr}, \lean{h.left}/\lean{h.right}, \lean{h ▸ e}
(transportar \lean{e} por la igualdad \lean{h}), \lean{rfl} (reflexividad) y
\lean{Termino.noConfusion} (constructores distintos producen valores
distintos).

%=====================================================================
\section{Catálogo de esquemas de demostración}\label{sec:esquemas}
%=====================================================================

Cada esquema tiene la misma estructura: la forma matemática, una plantilla
Lean genérica (en gris, escrita para este documento), la lectura en
español y uno o más ejemplos tomados literalmente de la tesis.

%---------------------------------------------------------------------
\subsection{Prueba directa: introducir \texorpdfstring{\(\forall\)}{∀} y \texorpdfstring{\(\to\)}{→}}\label{esq:directa}
\begin{esquema}{Esquema 1. Prueba directa}
\forma Para probar \(\forall x.\ P(x) \to Q(x)\): sea \(x\) arbitrario,
supongamos \(P(x)\) y deduzcamos \(Q(x)\).
\begin{codigolean}{Plantilla}
\begin{Verbatim}[commandchars=\\\{\}]
\LKw{theorem} t : ∀ x, P x → Q x := \LKw{by}
  \LTac{intro} x hP          \LCom{-- sea x; supongamos P x}
  \LTac{exact} …             \LCom{-- un término de tipo Q x}
\end{Verbatim}
\end{codigolean}
\lectura \lean{intro} nombra el objeto arbitrario y la suposición; \lean{exact}
entrega la conclusión. En Lean, \(\neg P\) es \(P \to \mathrm{False}\), así que
este esquema también prueba negaciones.
\donde En todas partes. Ejemplos mínimos: \lean{identidad_es_tautologia},
\lean{tautologia_de_derivable}, \lean{reduceBeta_reflexiva}.
\end{esquema}
\begin{codigolean}{\ruta{Thesis/Prop/Main.lean}{11}}
\begin{Verbatim}[commandchars=\\\{\}]
\LKw{theorem} identidad_es_tautologia : EsTautologia tautologiaIdentidad := \LKw{by}
  \LCom{-- A probar: P → P bajo toda valuación. Método: introducción de →.}
  \LTac{intro} v
  \LTac{change} evaluar v P → evaluar v P
  \LTac{intro} hP
  \LTac{exact} hP
\end{Verbatim}
\end{codigolean}
\lean{EsTautologia φ} se define como \lean{∀ v : Valuacion, evaluar v φ}; por
eso \lean{intro v} es posible. \lean{change} hace visible que, para
\lean{P ⟶ P}, \lean{evaluar} produce una implicación de Lean.

%---------------------------------------------------------------------
\subsection{Razonamiento hacia adelante con \texttt{have}}\label{esq:adelante}
\begin{esquema}{Esquema 2. Cadena de hechos intermedios}
\forma «De \(A\) obtenemos \(B\); de \(B\) y \(C\), \(D\); luego \(E\)».
\begin{codigolean}{Plantilla}
\begin{Verbatim}[commandchars=\\\{\}]
  \LKw{have} hB : B := f hA
  \LKw{have} hD : D := g hB hC
  \LTac{exact} k hD
\end{Verbatim}
\end{codigolean}
\lectura Cada \lean{have} es una línea de la prueba en papel con su
justificación a la derecha de \lean{:=}. El tipo escrito es la afirmación; Lean
comprueba que la justificación tiene exactamente ese tipo.
\donde Es el estilo dominante de la tesis: \lean{introduccion_doble_negacion},
\lean{analisis_por_casos}, \lean{completitud_ND}, \lean{quicksort_correcto},
\lean{unicidad_forma_normal}.
\end{esquema}
\begin{codigolean}{\ruta{Thesis/Prop/NaturalDeduction.lean}{96}}
\begin{Verbatim}[commandchars=\\\{\}]
\LKw{theorem} introduccion_doble_negacion \{Γ φ\} (d : ND Γ φ) : ND Γ (~~φ) := \LKw{by}
  \LCom{-- A probar: ND Γ (~~φ). Método: introducción de negación, no RAA.}
  \LKw{let} Δ := insert (~φ) Γ
  \LKw{have} hNeg : (~φ) ∈ Δ := Set.mem_insert _ _
  \LKw{have} hInclusion : Γ ⊆ Δ := Set.subset_insert _ _
  \LKw{have} dNeg : ND Δ (~φ) := ND.hyp hNeg
  \LKw{have} dPos : ND Δ φ := debilitamiento d hInclusion
  \LKw{have} dFalsedad : ND Δ Falsedad := ND.negE dNeg dPos
  \LCom{-- negI descarga el supuesto ~φ, no la premisa d.}
  \LTac{exact} ND.negI dFalsedad
\end{Verbatim}
\end{codigolean}
En papel: «Sea \(\Delta = \Gamma \cup \{\neg\varphi\}\). En \(\Delta\) tenemos
\(\neg\varphi\) por hipótesis y \(\varphi\) por debilitamiento de \(d\); por
\(\neg\)E obtenemos \(\bot\), y por \(\neg\)I, descargando \(\neg\varphi\),
\(\Gamma \vdash \neg\neg\varphi\)». Lean comprueba cada línea; en particular
comprueba que \lean{ND.negI} descarga justamente \lean{~φ}, porque su premisa
tiene tipo \lean{ND (insert (~φ) Γ) Falsedad}.

%---------------------------------------------------------------------
\subsection{Bicondicionales}\label{esq:iff}
\begin{esquema}{Esquema 3. Probar y usar \texorpdfstring{\(\leftrightarrow\)}{↔}}
\forma \(P \leftrightarrow Q\) se prueba con dos implicaciones. Para usarlo,
se toma una dirección o se reescribe con él.
\begin{codigolean}{Plantilla}
\begin{Verbatim}[commandchars=\\\{\}]
  \LTac{constructor}
  · \LTac{intro} hP; \LTac{exact} …   \LCom{-- (⇒)}
  · \LTac{intro} hQ; \LTac{exact} …   \LCom{-- (⇐)}
\LCom{-- uso:  h.mp hP : Q     h.mpr hQ : P     rw [h]}
\end{Verbatim}
\end{codigolean}
\donde \lean{correccion_completitud}, \lean{evaluar_falsedad},
\lean{relacion_pares_cons_iff}, \lean{evaluar_traduccion},
\lean{alfaEq_iff_alfaTesis}.
\end{esquema}
\begin{codigolean}{\ruta{Thesis/Prop/Completeness.lean}{275}}
\begin{Verbatim}[commandchars=\\\{\}]
\LKw{theorem} correccion_completitud (φ : Formula) : EsTautologia φ ↔ ND (∅ : Set Formula) φ
    \LCont{↪ }:= \LKw{by}
  \LCom{-- A probar: ambas direcciones. Método: introducción del bicondicional.}
  \LTac{constructor}
  · \LTac{intro} hTautologia
    \LTac{exact} completitud_ND φ hTautologia
  · \LTac{intro} d
    \LTac{exact} tautologia_de_derivable d
\end{Verbatim}
\end{codigolean}
El teorema final del capítulo~11 es la conjunción de los dos grandes resultados:
completitud (\(\Rightarrow\)) y corrección (\(\Leftarrow\)). Las dos metas que
produce \lean{constructor} son:
\begin{estadolean}{Estado después de \texttt{constructor}}
\begin{Verbatim}[commandchars=\\\{\}]
\LCom{-- meta 1: case mp}
φ : Formula
⊢ EsTautologia φ → ND ∅ φ

\LCom{-- meta 2: case mpr}
φ : Formula
⊢ ND ∅ φ → EsTautologia φ
\end{Verbatim}
\end{estadolean}

%---------------------------------------------------------------------
\subsection{Conjunciones}\label{esq:conj}
\begin{esquema}{Esquema 4. Introducir y eliminar \texorpdfstring{\(\wedge\)}{∧}}
\forma Para probar \(P \wedge Q\) se prueban \(P\) y \(Q\). De
\(P \wedge Q\) se obtienen \(P\) y \(Q\).
\begin{codigolean}{Plantilla}
\begin{Verbatim}[commandchars=\\\{\}]
  \LTac{refine} ⟨?_, ?_⟩            \LCom{-- dos metas: P y Q}
  \LCom{-- o bien:  exact ⟨hP, hQ⟩    exact And.intro hP hQ}
  \LTac{obtain} ⟨hP, hQ⟩ := h        \LCom{-- uso; también h.left, h.right}
\end{Verbatim}
\end{codigolean}
\donde \lean{kalmar} (cuyo enunciado es una conjunción, la parte positiva y
la parte negativa), \lean{quicksort_correcto}, \lean{descarga}
(\lean{hPartes.left}), \lean{lema_7_6}, \lean{lema_7_31}, \lean{lema_7_32}.
\end{esquema}
\begin{codigolean}{\ruta{Thesis/Sort/Quicksort.lean}{322}}
\begin{Verbatim}[commandchars=\\\{\}]
\LKw{theorem} quicksort_correcto (l : List α) :
    quicksort l ~ l ∧ Ordenada (quicksort l) := \LKw{by}
  \LCom{-- A probar: permutación ∧ ordenamiento para la misma lista l.}
  \LCom{-- Método: introducción de la conjunción a partir de los dos teoremas.}
  \LKw{have} hPermutacion : quicksort l ~ l := quicksort_permutacion l
  \LKw{have} hOrdenada : Ordenada (quicksort l) := quicksort_ordenada l
  \LTac{exact} And.intro hPermutacion hOrdenada
\end{Verbatim}
\end{codigolean}
Una conjunción también puede ser la hipótesis de una regla. En
\lean{relacion_pares_cons_iff} la dirección (\(\Rightarrow\)) \emph{elimina}
el constructor \lean{Pairwise.cons} con \lean{cases}, y la dirección
(\(\Leftarrow\)) lo \emph{introduce}:
\begin{codigolean}{\ruta{Thesis/Sort/Quicksort.lean}{73}}
\begin{Verbatim}[commandchars=\\\{\}]
\LKw{theorem} relacion_pares_cons_iff (R : α → α → Prop) (a : α) (l : List α) :
    List.Pairwise R (a :: l) ↔ (∀ b ∈ l, R a b) ∧ List.Pairwise R l := \LKw{by}
  \LCom{-- A probar: equivalencia entre Pairwise y las dos premisas del constructor cons.}
  \LCom{-- Por eliminación del constructor en una dirección e introducción en la otra.}
  \LTac{constructor}
  · \LTac{intro} h
    \LCom{-- Se extraen la relación de la cabeza con la cola y Pairwise de la cola.}
    \LTac{cases} h \LKw{with}
    | cons hCabeza hCola => \LTac{exact} ⟨hCabeza, hCola⟩
  · \LTac{rintro} ⟨hCabeza, hCola⟩
    \LCom{-- Las mismas premisas reconstruyen Pairwise para a :: l.}
    \LTac{exact} List.Pairwise.cons hCabeza hCola
\end{Verbatim}
\end{codigolean}

%---------------------------------------------------------------------
\subsection{Disyunciones y análisis por casos de una hipótesis}\label{esq:disy}
\begin{esquema}{Esquema 5. Eliminar \texorpdfstring{\(\vee\)}{∨}: prueba por casos}
\forma Si \(A \vee B\), y \(C\) se sigue de \(A\) y también de \(B\),
entonces \(C\).
\begin{codigolean}{Plantilla}
\begin{Verbatim}[commandchars=\\\{\}]
  \LTac{rcases} h \LKw{with} hA | hB
  · …   \LCom{-- caso A, con hA : A}
  · …   \LCom{-- caso B, con hB : B}
\end{Verbatim}
\end{codigolean}
\lectura \lean{rcases} crea una meta por cada disyunto, con la misma
conclusión. Para \emph{probar} una disyunción se elige un lado:
\lean{Or.inl} o \lean{Or.inr}.
\donde \lean{satisfaccion_insertar}, \lean{pertenencia_literal},
\lean{sub_transitiva}, \lean{igualBeta_de_reduceBeta},
\lean{church_rosser_conversion}, \lean{quicksort_ordenada} (condición
cruzada), \lean{variables_libres}.
\end{esquema}
\begin{codigolean}{\ruta{Thesis/Prop/NaturalDeduction.lean}{146}}
\begin{Verbatim}[commandchars=\\\{\}]
\LKw{theorem} satisfaccion_insertar \{v : Valuacion\} \{Γ : Set Formula\} \{χ : Formula\}
    (hχ : evaluar v χ) (hΓ : ∀ ψ ∈ Γ, evaluar v ψ) :
    ∀ ψ ∈ insert χ Γ, evaluar v ψ := \LKw{by}
  \LCom{-- A probar: cada fórmula del contexto extendido es verdadera bajo v.}
  \LCom{-- Método: casos de pertenencia; ψ = χ o ψ ∈ Γ.}
  \LTac{intro} ψ hψ
  \LKw{have} hCasos : ψ = χ ∨ ψ ∈ Γ := Set.mem_insert_iff.mp hψ
  \LTac{rcases} hCasos \LKw{with} hIgualdad | hEnGamma
  · \LTac{subst} ψ
    \LTac{exact} hχ
  · \LTac{exact} hΓ ψ hEnGamma
\end{Verbatim}
\end{codigolean}
La pertenencia a \lean{insert χ Γ} se reescribe como disyunción
(\lean{Set.mem_insert_iff}), y \lean{rcases} la divide en dos ramas:
\begin{estadolean}{Estado antes}
\begin{Verbatim}[commandchars=\\\{\}]
v : Valuacion
Γ : Set Formula
χ : Formula
hχ : evaluar v χ
hΓ : ∀ ψ ∈ Γ, evaluar v ψ
ψ : Formula
hψ : ψ ∈ insert χ Γ
hCasos : ψ = χ ∨ ψ ∈ Γ
⊢ evaluar v ψ
\end{Verbatim}
\end{estadolean}
\tacticaentre{rcases hCasos with hIgualdad | hEnGamma}
\begin{estadolean}{Estado después}
\begin{Verbatim}[commandchars=\\\{\}]
\LCom{-- meta 1: case inl}
v : Valuacion
Γ : Set Formula
χ : Formula
hχ : evaluar v χ
hΓ : ∀ ψ ∈ Γ, evaluar v ψ
ψ : Formula
hψ : ψ ∈ insert χ Γ
hIgualdad : ψ = χ
⊢ evaluar v ψ

\LCom{-- meta 2: case inr}
v : Valuacion
Γ : Set Formula
χ : Formula
hχ : evaluar v χ
hΓ : ∀ ψ ∈ Γ, evaluar v ψ
ψ : Formula
hψ : ψ ∈ insert χ Γ
hEnGamma : ψ ∈ Γ
⊢ evaluar v ψ
\end{Verbatim}
\end{estadolean}
En la primera rama \lean{subst} reemplaza \(\psi\) por \(\chi\) y la meta queda
cerrada por \lean{hχ}; en la segunda basta \lean{hΓ}. Así se introduce una
disyunción, en la prueba de que la semántica de la tesis coincide con la de
Foundation:
\begin{codigolean}{\ruta{Thesis/Prop/BridgeLemma.lean}{29} (fragmento)}
\begin{Verbatim}[commandchars=\\\{\}]
· \LTac{intro} hEvaluacionImp
  \LCom{-- Tercero excluido metateórico para construir la disyunción.}
  \LTac{by_cases} hP : evaluar v p
  · \LKw{have} hQ : evaluar v q := hEvaluacionImp hP
    \LTac{exact} Or.inr (hipIndQ.mp hQ)
  · \LKw{have} hNoSatisfaccion : ¬ Satisface v (traducir p) := \LKw{by}
      \LTac{intro} hSatisfaccion
      \LTac{exact} hP (hipIndP.mpr hSatisfaccion)
    \LTac{exact} Or.inl hNoSatisfaccion
\end{Verbatim}
\end{codigolean}

%---------------------------------------------------------------------
\subsection{Casos sobre una proposición: \texttt{by\_cases}}\label{esq:bycases}
\begin{esquema}{Esquema 6. Tercero excluido en la metalógica}
\forma «O bien \(P\), o bien \(\neg P\). Si \(P\), \ldots; si no, \ldots».
\begin{codigolean}{Plantilla}
\begin{Verbatim}[commandchars=\\\{\}]
  \LTac{by_cases} hp : P
  · …   \LCom{-- con hp : P}
  · …   \LCom{-- con hp : ¬ P}
\end{Verbatim}
\end{codigolean}
\lectura A diferencia de \lean{rcases}, aquí no hay una disyunción previa.
Si \(P\) es decidible (por ejemplo, la igualdad de dos cadenas \lean{z = x}),
Lean usa su procedimiento de decisión. Si no lo es (por ejemplo,
\lean{evaluar v p}, que es una proposición arbitraria), usa el tercero
excluido clásico, que depende del axioma \lean{Classical.choice}; por eso la
completitud aparece con ese axioma en la auditoría.
\donde \lean{congruencia_literal}, \lean{kalmar} (caso \(p \to q\)),
\lean{evaluar_traduccion}, \lean{sustLibre_conmuta} y \lean{lema_sustitucion}
(los tres subcasos de la variable: \(z \equiv x\), \(z \equiv y\), otro).
\end{esquema}
\begin{codigolean}{\ruta{Thesis/Prop/Completeness.lean}{155}}
\begin{Verbatim}[commandchars=\\\{\}]
\LKw{theorem} congruencia_literal \{w1 w2 : Valuacion\} \{x : String\} (h : w1 x ↔ w2 x) :
    literal w1 x = literal w2 x := \LKw{by}
  \LCom{-- A probar: igualdad de literales. Método: casos sobre w1 x.}
  \LTac{by_cases} hx : w1 x
  · \LKw{have} hx2 : w2 x := h.mp hx
    \LTac{rw} [literal, literal, if_pos hx, if_pos hx2]
  · \LKw{have} hx2 : ¬ w2 x := \LKw{by}
      \LTac{intro} hVerdadero
      \LKw{have} hx1 : w1 x := h.mpr hVerdadero
      \LTac{exact} hx hx1
    \LTac{rw} [literal, literal, if_neg hx, if_neg hx2]
\end{Verbatim}
\end{codigolean}
En el Lema de Kalmár la división es semántica: se distingue según el valor de
verdad de \(p\) bajo \(v\), no se usa el lema objeto \lean{analisis_por_casos}:
\begin{estadolean}{Estado antes}
\begin{Verbatim}[commandchars=\\\{\}]
\LCom{-- (… 10 hipótesis omitidas …)}
hipIndPNeg : ¬evaluar v p → ND (contextoLiterales v listaAtomos) (~p)
hipIndQPos : evaluar v q → ND (contextoLiterales v listaAtomos) q
\LCom{-- (… 1 hipótesis omitidas …)}
hev : evaluar v (p ⟶ q)
⊢ ND (contextoLiterales v listaAtomos) (p ⟶ q)
\end{Verbatim}
\end{estadolean}
\tacticaentre{by\_cases hp : evaluar v p}
\begin{estadolean}{Estado después}
\begin{Verbatim}[commandchars=\\\{\}]
\LCom{-- meta 1: case pos}
\LCom{-- (… 10 hipótesis omitidas …)}
hipIndPNeg : ¬evaluar v p → ND (contextoLiterales v listaAtomos) (~p)
hipIndQPos : evaluar v q → ND (contextoLiterales v listaAtomos) q
\LCom{-- (… 1 hipótesis omitidas …)}
hev : evaluar v (p ⟶ q)
hp : evaluar v p
⊢ ND (contextoLiterales v listaAtomos) (p ⟶ q)

\LCom{-- meta 2: case neg}
\LCom{-- (… 10 hipótesis omitidas …)}
hipIndPNeg : ¬evaluar v p → ND (contextoLiterales v listaAtomos) (~p)
hipIndQPos : evaluar v q → ND (contextoLiterales v listaAtomos) q
\LCom{-- (… 1 hipótesis omitidas …)}
hev : evaluar v (p ⟶ q)
hp : ¬evaluar v p
⊢ ND (contextoLiterales v listaAtomos) (p ⟶ q)
\end{Verbatim}
\end{estadolean}

%---------------------------------------------------------------------
\subsection{Negación, contradicción y reducción al absurdo}\label{esq:negacion}
\begin{esquema}{Esquema 7. Probar \texorpdfstring{\(\neg P\)}{¬P}, usar \texttt{False}, razonar por el absurdo}
\forma (a) Para \(\neg P\): suponer \(P\) y llegar a una contradicción.
(b) De una contradicción se sigue cualquier cosa. (c) Para probar \(P\)
(clásicamente): suponer \(\neg P\) y llegar a una contradicción.
\begin{codigolean}{Plantilla}
\begin{Verbatim}[commandchars=\\\{\}]
  \LTac{intro} hP; …; \LTac{exact} hNoQ hQ       \LCom{-- (a)  ¬P = P → False}
  \LTac{exact} False.elim hFalso          \LCom{-- (b)  también: absurd hQ hNoQ}
  \LTac{by_contra} hNoP; …; \LTac{exact} …       \LCom{-- (c)  la meta pasa a ser False}
\end{Verbatim}
\end{codigolean}
\lectura Una negación \lean{hNoQ : ¬ Q} es una función: aplicada a
\lean{hQ : Q} produce \lean{False}. \lean{by_contra} y
\lean{Classical.byContradiction} son la reducción al absurdo \emph{de la
metalógica}.
\donde (a) \lean{no_derivable_falsedad}, \lean{formaNormal_sin_paso};
(b) \lean{pertenencia_literal}, \lean{variables_libres}, \lean{paso_de_redex};
(c) \lean{correccion} (caso \lean{classical}), \lean{kalmar}.
\end{esquema}
\begin{codigolean}{\ruta{Thesis/Prop/NaturalDeduction.lean}{218}}
\begin{Verbatim}[commandchars=\\\{\}]
\LKw{theorem} no_derivable_falsedad : ¬ ND (∅ : Set Formula) Falsedad := \LKw{by}
  \LCom{-- A probar: la inexistencia de una derivación cerrada de Falsedad.}
  \LCom{-- Método: introducción de negación; una tautología no puede ser siempre falsa.}
  \LTac{intro} d
  \LKw{let} v : Valuacion := \LKw{fun} _ => True
  \LKw{have} hTautologia : EsTautologia Falsedad := tautologia_de_derivable d
  \LKw{have} hFalsedad : evaluar v Falsedad := hTautologia v
  \LTac{exact} (evaluar_falsedad v).mp hFalsedad
\end{Verbatim}
\end{codigolean}
Consistencia: si hubiera una derivación cerrada de \(\bot\), por corrección
sería una tautología, y evaluada en la valuación constante verdadera daría
\lean{False}. Un ejemplo de (c), dentro de Kalmár: de
\(\neg(v(p) \to v(q))\) se deduce \(v(p)\) por el absurdo:
\begin{codigolean}{\ruta{Thesis/Prop/Completeness.lean}{121} (fragmento)}
\begin{Verbatim}[commandchars=\\\{\}]
\LKw{have} hp : evaluar v p := \LKw{by}
  \LCom{-- Por RAA metateórica, se supone que evaluar v p es falsa.}
  \LTac{by_contra} hNoP
  \LKw{have} hImp : evaluar v p → evaluar v q := \LKw{by}
    \LTac{intro} hP
    \LKw{have} hFalso : False := hNoP hP
    \LTac{exact} False.elim hFalso
  \LTac{exact} hNoImp hImp
\end{Verbatim}
\end{codigolean}
Para descartar casos imposibles sobre la forma de un término se usa que
constructores distintos dan valores distintos (\lean{noConfusion}); así ocurre
en \lean{paso_de_redex}, donde una variable o una abstracción no pueden ser un
redex.

%---------------------------------------------------------------------
\subsection{Existenciales: dar un testigo y usarlo}\label{esq:existe}
\begin{esquema}{Esquema 8. Introducir y eliminar \texorpdfstring{\(\exists\)}{∃}}
\forma Para probar \(\exists x.\ P(x)\) se exhibe un \(a\) con \(P(a)\). De
\(\exists x.\ P(x)\) se puede «tomar un \(x\) tal que \(P(x)\)».
\begin{codigolean}{Plantilla}
\begin{Verbatim}[commandchars=\\\{\}]
  \LTac{obtain} ⟨x, hx⟩ := h          \LCom{-- uso: sea x con hx : P x}
  \LTac{exact} ⟨a, prueba_de_P_a⟩     \LCom{-- introducción: testigo y prueba}
\end{Verbatim}
\end{codigolean}
\donde Las relaciones \(\twoheadrightarrow_\beta\) y \(=_\beta\) se definen como
«existe \(n\) y una cadena de \(n\) pasos», de modo que sus propiedades se
prueban con este esquema: \lean{reduceBeta_transitiva},
\lean{igualBeta_simetrica}, \lean{church_rosser_conversion},
\lean{paso_de_redex}. También \lean{dominio_extender}.
\end{esquema}
\begin{codigolean}{\ruta{Thesis/Teoria/Lambda.lean}{185}}
\begin{Verbatim}[commandchars=\\\{\}]
\LKw{def} ReduceBeta (M N : Termino) : Prop :=
  ∃ n : Nat, Cadena (PasoBeta sust) n M N
\end{Verbatim}
\end{codigolean}
\begin{codigolean}{\ruta{Thesis/Teoria/Lambda.lean}{213}}
\begin{Verbatim}[commandchars=\\\{\}]
\LKw{theorem} reduceBeta_transitiva \{L M N : Termino\}
    (hLM : ReduceBeta sust L M) (hMN : ReduceBeta sust M N) : ReduceBeta sust L N := \LKw{by}
  \LCom{-- Secuencia de `r ≥ 0` reducciones de `L` a `M`.}
  \LTac{obtain} ⟨r, hCadenaLM⟩ := hLM
  \LCom{-- Secuencia de `s ≥ 0` reducciones de `M` a `N`.}
  \LTac{obtain} ⟨s, hCadenaMN⟩ := hMN
  \LCom{-- Se concatenan: una secuencia de `r + s` pasos de `L` a `N`.}
  \LTac{exact} ⟨r + s, Cadena.concatenar hCadenaLM hCadenaMN⟩
\end{Verbatim}
\end{codigolean}
En papel: «tomemos una secuencia de \(r\) pasos de \(L\) a \(M\) y una de \(s\)
pasos de \(M\) a \(N\); al concatenarlas se obtiene una de \(r+s\) pasos de
\(L\) a \(N\)». El testigo es \lean{r + s} y la prueba es
\lean{Cadena.concatenar}.

%---------------------------------------------------------------------
\subsection{Igualdades: reescritura, sustitución y cadenas}\label{esq:igualdad}
\begin{esquema}{Esquema 9. Razonamiento ecuacional}
\forma «Sustituyendo \(a\) por \(b\) se obtiene\ldots»; «\(a = b = c\)».
\begin{codigolean}{Plantilla}
\begin{Verbatim}[commandchars=\\\{\}]
  \LTac{rw} [eq₁, eq₂]            \LCom{-- reescribe la meta de izquierda a derecha}
  \LTac{rw} [← eq₁] \LKw{at} h          \LCom{-- de derecha a izquierda, en la hipótesis h}
  \LTac{subst} hx                 \LCom{-- hx : x = a; reemplaza x por a en todo}
  \LTac{calc} a = b := p₁
    _ = c := p₂
\end{Verbatim}
\end{codigolean}
\lectura \lean{rw} usa una igualdad (o una equivalencia) como regla de
reescritura. También despliega definiciones dadas por ecuaciones:
\lean{rw [literal, if_pos hvs]} significa «por definición de \lean{literal}, y
como \(v(s)\) es verdadero, la condición toma la rama \emph{entonces}». Un
\lean{calc} reproduce literalmente una cadena de igualdades del papel; también
admite relaciones transitivas, como la permutación \lean{~}.
\donde \lean{congruencia_contextoLiterales}, \lean{kalmar} (caso atómico),
\lean{descarga}, \lean{quicksort_permutacion} (un \lean{calc} de
permutaciones), \lean{lema_sustitucion}, \lean{lema_sustitucion_alfa} (un
\lean{calc} de igualdades), \lean{formaNormal_reduce_igual}.
\end{esquema}
\begin{codigolean}{\ruta{Thesis/Prop/Completeness.lean}{59} (fragmento)}
\begin{Verbatim}[commandchars=\\\{\}]
| atom s =>
    \LTac{intro} hsub
    \LKw{have} hAtomo : s ∈ atomos (.atom s) := List.mem_singleton_self s
    \LKw{have} hs : s ∈ listaAtomos := hsub hAtomo
    \LTac{refine} ⟨?_, ?_⟩
    · \LTac{intro} hev
      \LKw{have} hvs : v s := hev
      \LKw{have} hmem : literal v s ∈ contextoLiterales v listaAtomos := pertenencia_literal
          \LCont{↪ }(v := v) hs
      \LTac{rw} [literal, if_pos hvs] \LKw{at} hmem
      \LTac{exact} ND.hyp hmem
    · \LTac{intro} hev
      \LKw{have} hvs : ¬ v s := hev
      \LKw{have} hmem : literal v s ∈ contextoLiterales v listaAtomos := pertenencia_literal
          \LCont{↪ }(v := v) hs
      \LTac{rw} [literal, if_neg hvs] \LKw{at} hmem
      \LTac{exact} ND.hyp hmem
\end{Verbatim}
\end{codigolean}
Efecto de la reescritura en la hipótesis \lean{hmem}:
\begin{estadolean}{Estado antes}
\begin{Verbatim}[commandchars=\\\{\}]
\LCom{-- (… 5 hipótesis omitidas …)}
hs : s ∈ listaAtomos
\LCom{-- (… 1 hipótesis omitidas …)}
hvs : v s
hmem : literal v s ∈ contextoLiterales v listaAtomos
⊢ ND (contextoLiterales v listaAtomos) (Formula.atom s)
\end{Verbatim}
\end{estadolean}
\tacticaentre{rw [literal, if\_pos hvs] at hmem}
\begin{estadolean}{Estado después}
\begin{Verbatim}[commandchars=\\\{\}]
\LCom{-- (… 5 hipótesis omitidas …)}
hs : s ∈ listaAtomos
\LCom{-- (… 1 hipótesis omitidas …)}
hvs : v s
hmem : Formula.atom s ∈ contextoLiterales v listaAtomos
⊢ ND (contextoLiterales v listaAtomos) (Formula.atom s)
\end{Verbatim}
\end{estadolean}

%---------------------------------------------------------------------
\subsection{Inducción estructural sobre datos}\label{esq:induccion-estructural}
\begin{esquema}{Esquema 10. Inducción sobre fórmulas, términos y listas}
\forma Para probar \(\forall \varphi.\ P(\varphi)\): probar \(P\) para cada
constructor, suponiendo \(P\) para los subcomponentes (hipótesis inductiva).
\begin{codigolean}{Plantilla}
\begin{Verbatim}[commandchars=\\\{\}]
\LKw{theorem} t : ∀ φ : Formula, P φ := \LKw{by}
  \LTac{intro} φ
  \LTac{induction} φ \LKw{with}
  | atom s => …                          \LCom{-- caso base}
  | neg p hipInd => …                    \LCom{-- hipInd : P p}
  | impl p q hipIndP hipIndQ => …        \LCom{-- dos hipótesis}
\end{Verbatim}
\end{codigolean}
\lectura Lean genera una meta por constructor y nombra las hipótesis
inductivas según el patrón. La regla de oro: \textbf{lo que deba variar
en la hipótesis inductiva debe seguir en la meta cuando se aplica
\lean{induction}}. Por eso Kalmár deja la premisa
\lean{atomos φ ⊆ listaAtomos → …} dentro del enunciado y hace \lean{intro hsub}
\emph{después} de la inducción: así la hipótesis inductiva es una implicación
que se puede aplicar a la cobertura de \(p\) o de \(q\).
\donde Fórmulas: \lean{kalmar}, \lean{evaluar_traduccion}. Listas:
\lean{pertenencia_literal}, \lean{congruencia_contextoLiterales},
\lean{descarga}, \lean{longitud_filtro_le}, \lean{particion_filtros_permutacion},
\lean{relacion_pares_concatenacion_iff}. Términos~\(\lambda\):
\lean{sub_transitiva}, \lean{paso_de_redex}, \lean{lema_sustitucion},
\lean{sust_no_libre}, \lean{sustLibre_conmuta}.
\end{esquema}
\begin{codigolean}{\ruta{Thesis/Prop/Completeness.lean}{32}}
\begin{Verbatim}[commandchars=\\\{\}]
\LKw{theorem} pertenencia_literal \{v : Valuacion\} \{p : String\} \{listaAtomos : List String\}
    (h : p ∈ listaAtomos) : literal v p ∈ contextoLiterales v listaAtomos := \LKw{by}
  \LCom{-- A probar: pertenencia del literal correspondiente al átomo p.}
  \LCom{-- Método: inducción en listaAtomos; la premisa p ∈ listaAtomos se conserva en la}
      \LCont{↪ }\LCom{hipótesis inductiva.}
  \LTac{induction} listaAtomos \LKw{with}
  | nil =>
      \LKw{have} hFalso : False := List.not_mem_nil h
      \LTac{exact} False.elim hFalso
  | cons a cola hipInd =>
      \LKw{have} hCasos : p = a ∨ p ∈ cola := List.mem_cons.mp h
      \LTac{rcases} hCasos \LKw{with} hIgualdad | hCola
      · \LTac{subst} p
        \LTac{exact} Set.mem_insert _ _
      · \LKw{have} hLiteralCola : literal v p ∈ contextoLiterales v cola := hipInd hCola
        \LTac{exact} Set.mem_insert_of_mem (literal v a) hLiteralCola
\end{Verbatim}
\end{codigolean}
Al dividir la meta de Kalmár en sus dos partes, en el caso \lean{neg} ya se
dispone de las dos mitades de la hipótesis inductiva:
\begin{estadolean}{Estado antes}
\begin{Verbatim}[commandchars=\\\{\}]
\LCom{-- (… 4 hipótesis omitidas …)}
hsub : atomos (~p) ⊆ listaAtomos
\LCom{-- (… 1 hipótesis omitidas …)}
hipIndPos : evaluar v p → ND (contextoLiterales v listaAtomos) p
hipIndNeg : ¬evaluar v p → ND (contextoLiterales v listaAtomos) (~p)
⊢ (evaluar v (~p) → ND (contextoLiterales v listaAtomos) (~p)) ∧
  (¬evaluar v (~p) → ND (contextoLiterales v listaAtomos) (~~p))
\end{Verbatim}
\end{estadolean}
\tacticaentre{refine ⟨?\_, ?\_⟩}
\begin{estadolean}{Estado después}
\begin{Verbatim}[commandchars=\\\{\}]
\LCom{-- meta 1: case neg.refine_1}
\LCom{-- (… 4 hipótesis omitidas …)}
hsub : atomos (~p) ⊆ listaAtomos
\LCom{-- (… 1 hipótesis omitidas …)}
hipIndPos : evaluar v p → ND (contextoLiterales v listaAtomos) p
hipIndNeg : ¬evaluar v p → ND (contextoLiterales v listaAtomos) (~p)
⊢ evaluar v (~p) → ND (contextoLiterales v listaAtomos) (~p)

\LCom{-- meta 2: case neg.refine_2}
\LCom{-- (… 4 hipótesis omitidas …)}
hsub : atomos (~p) ⊆ listaAtomos
\LCom{-- (… 1 hipótesis omitidas …)}
hipIndPos : evaluar v p → ND (contextoLiterales v listaAtomos) p
hipIndNeg : ¬evaluar v p → ND (contextoLiterales v listaAtomos) (~p)
⊢ ¬evaluar v (~p) → ND (contextoLiterales v listaAtomos) (~~p)
\end{Verbatim}
\end{estadolean}

\begin{nota}[El caso de \texttt{descarga}]
En \lean{descarga} la hipótesis inductiva debe valer \emph{para toda
valuación}, porque se aplica a valuaciones modificadas con
\lean{Function.update}. Por eso el enunciado cuantifica
\lean{∀ v, ND (contextoLiterales v listaAtomos) φ} como premisa dentro de la
meta, y la prueba construye primero \lean{hUniversal} y solo al final aplica
\lean{hipInd}. Es el mismo principio: generalizar antes de inducir.
\end{nota}

%---------------------------------------------------------------------
\subsection{Inducción sobre derivaciones (reglas)}\label{esq:induccion-derivaciones}
\begin{esquema}{Esquema 11. Inducción sobre un predicado inductivo}
\forma «Por inducción sobre la derivación de \(\Gamma \vdash \varphi\)»: se
prueba que cada regla preserva la propiedad.
\begin{codigolean}{Plantilla}
\begin{Verbatim}[commandchars=\\\{\}]
\LKw{theorem} t \{Γ φ\} (d : ND Γ φ) : ∀ \{Δ\}, Γ ⊆ Δ → ND Δ φ := \LKw{by}
  \LTac{induction} d \LKw{with}
  | @hyp Γ φ hmem => …               \LCom{-- una rama por regla}
  | @impI Γ φ ψ _ hipInd => …        \LCom{-- hipInd: propiedad de la premisa}
  …
\end{Verbatim}
\end{codigolean}
\lectura Cada rama corresponde a la \emph{última regla} usada en la
derivación. Las hipótesis inductivas dicen que las premisas de esa regla ya
cumplen la propiedad. El \lean{@} permite nombrar los argumentos implícitos de
la regla (aquí \lean{Γ}, \lean{φ}, \lean{ψ}). Como en el esquema anterior, lo
que deba variar (el contexto mayor \(\Delta\)) se deja cuantificado en la
meta.
\donde \lean{debilitamiento} (capítulos 9 y 11), \lean{correccion},
\lean{variables_libres}, \lean{Deriva.bienFormado}, \lean{formaNormal_sin_paso},
\lean{alfaTesis_implica_alfaEq} y todos los lemas de \lean{Cadena}
(\lean{concatenar}, \lean{anteponer}, \lean{cambiar}, \lean{invertir},
\lean{formaNormal_reduce_igual}).
\end{esquema}
\begin{codigolean}{\ruta{Thesis/Prop/NaturalDeduction.lean}{51}}
\begin{Verbatim}[commandchars=\\\{\}]
\LKw{theorem} debilitamiento \{Γ φ\} (d : ND Γ φ) : ∀ \{Δ\}, Γ ⊆ Δ → ND Δ φ := \LKw{by}
  \LCom{-- A probar: la misma conclusión bajo cualquier contexto mayor Δ.}
  \LCom{-- Método: inducción sobre d; Δ permanece cuantificado en cada hipótesis inductiva.}
  \LTac{induction} d \LKw{with}
  | @hyp Γ φ hmem =>
      \LTac{intro} Δ hsub
      \LKw{have} hEnDelta : φ ∈ Δ := hsub hmem
      \LTac{exact} ND.hyp hEnDelta
  | @impI Γ φ ψ _ hipInd =>
      \LTac{intro} Δ hsub
      \LKw{have} hAmpliado : insert φ Γ ⊆ insert φ Δ := Set.insert_subset_insert hsub
      \LKw{have} dCuerpo : ND (insert φ Δ) ψ := hipInd hAmpliado
      \LCom{-- impI descarga φ; las fórmulas de Δ permanecen.}
      \LTac{exact} ND.impI dCuerpo
  | @impE Γ φ ψ _ _ hipIndImp hipIndAntecedente =>
      \LTac{intro} Δ hsub
      \LKw{have} dImp : ND Δ (φ ⟶ ψ) := hipIndImp hsub
      \LKw{have} dAntecedente : ND Δ φ := hipIndAntecedente hsub
      \LTac{exact} ND.impE dImp dAntecedente
  | @negI Γ φ _ hipInd =>
      \LTac{intro} Δ hsub
      \LKw{have} hAmpliado : insert φ Γ ⊆ insert φ Δ := Set.insert_subset_insert hsub
      \LKw{have} dFalsedad : ND (insert φ Δ) Falsedad := hipInd hAmpliado
      \LTac{exact} ND.negI dFalsedad
  | @negE Γ φ _ _ hipIndNeg hipIndPos =>
      \LTac{intro} Δ hsub
      \LKw{have} dNeg : ND Δ (~φ) := hipIndNeg hsub
      \LKw{have} dPos : ND Δ φ := hipIndPos hsub
      \LTac{exact} ND.negE dNeg dPos
  | @botE Γ φ _ hipInd =>
      \LTac{intro} Δ hsub
      \LKw{have} dFalsedad : ND Δ Falsedad := hipInd hsub
      \LTac{exact} ND.botE dFalsedad
  | @classical Γ φ _ hipInd =>
      \LTac{intro} Δ hsub
      \LKw{have} hAmpliado : insert (~φ) Γ ⊆ insert (~φ) Δ := Set.insert_subset_insert hsub
      \LKw{have} dFalsedad : ND (insert (~φ) Δ) Falsedad := hipInd hAmpliado
      \LCom{-- Por RAA objeto, se descarga ~φ y se concluye φ bajo Δ.}
      \LTac{exact} ND.classical dFalsedad
\end{Verbatim}
\end{codigolean}
La regla \lean{impI} es el caso que justifica el esquema: la premisa vive en el
contexto ampliado \lean{insert φ Γ}, así que la hipótesis inductiva debe
aplicarse a \lean{insert φ Γ ⊆ insert φ Δ}, no a \lean{Γ ⊆ Δ}. Si \(\Delta\)
hubiera quedado fijo, esa aplicación sería imposible. La lectura guiada de la
sección~\ref{sec:guiada-debilitamiento} recorre este caso estado por estado.

El mismo esquema, con las tres reglas del tipado simple, demuestra el Lema
7.57 (las variables libres de un término tipado están declaradas):
\begin{codigolean}{\ruta{Thesis/Teoria/TiposSimples.lean}{168}}
\begin{Verbatim}[commandchars=\\\{\}]
\LKw{theorem} variables_libres \{Γ : Contexto\} \{L : Termino\} \{σ : Tipo\}
    (d : Γ ⊢ L : σ) : variablesLibres L ⊆ dominio Γ := \LKw{by}
  \LCom{-- A probar: FV(L) ⊆ dom(Γ). Por inducción sobre la derivación de Γ ⊢ L : σ.}
  \LTac{induction} d \LKw{with}
  | @var Γ x σ _ hmem =>
      \LCom{-- Caso (variable): L ≡ x y la regla exige (x:σ) ∈ Γ.}
      \LCom{-- Por definición de dominio, x ∈ dom(Γ).}
      \LKw{have} hDominio : x ∈ dominio Γ := ⟨σ, hmem⟩
      \LCom{-- FV(x) = \{x\}.}
      \LKw{have} hFV : variablesLibres (.var x) = \{x\} := \LTac{rfl}
      \LTac{rw} [hFV]
      \LTac{exact} Set.singleton_subset_iff.mpr hDominio
  | @apl Γ M N σ τ _ _ hipIndM hipIndN =>
      \LCom{-- Caso (aplicación): L ≡ M N, con premisas Γ ⊢ M : σ → τ y Γ ⊢ N : σ.}
      \LCom{-- Hipótesis inductivas: FV(M) ⊆ dom(Γ) y FV(N) ⊆ dom(Γ).}
      \LCom{-- FV(M N) = FV(M) ∪ FV(N) ⊆ dom(Γ) ∪ dom(Γ) = dom(Γ).}
      \LKw{have} hFV : variablesLibres (.app M N) = variablesLibres M ∪ variablesLibres N :=
          \LCont{↪ }\LTac{rfl}
      \LTac{rw} [hFV]
      \LTac{exact} Set.union_subset hipIndM hipIndN
  | @abst Γ x σ τ M _ hipInd =>
      \LCom{-- Caso (abstracción): L ≡ λx:σ. M, con premisa Γ, x:σ ⊢ M : τ.}
      \LCom{-- Hipótesis inductiva: FV(M) ⊆ dom(Γ, x:σ) = dom(Γ) ∪ \{x\}.}
      \LKw{have} hCuerpo : variablesLibres M ⊆ dominio Γ ∪ \{x\} := \LKw{by}
        \LTac{rw} [← dominio_extender Γ x σ]
        \LTac{exact} hipInd
      \LCom{-- FV(λx:σ. M) = FV(M) \textbackslash{} \{x\} ⊆ (dom(Γ) ∪ \{x\}) \textbackslash{} \{x\} ⊆ dom(Γ).}
      \LKw{have} hFV : variablesLibres (.abs x σ M) = variablesLibres M \textbackslash{} \{x\} := \LTac{rfl}
      \LTac{rw} [hFV]
      \LTac{intro} y hy
      \LKw{have} hyCuerpo : y ∈ dominio Γ ∪ \{x\} := hCuerpo hy.1
      \LTac{rcases} hyCuerpo \LKw{with} hyDominio | hyx
      · \LTac{exact} hyDominio
      · \LTac{exact} absurd hyx hy.2
\end{Verbatim}
\end{codigolean}
\lean{Cadena} es un predicado inductivo con dos reglas (cadena vacía, y cadena
más un paso). La concatenación se prueba por inducción sobre la
\emph{segunda} cadena:
\begin{codigolean}{\ruta{Thesis/Teoria/Lambda.lean}{117}}
\begin{Verbatim}[commandchars=\\\{\}]
\LKw{inductive} Cadena \{α : Type\} (R : α → α → Prop) : Nat → α → α → Prop
  \LCom{/-- La cadena vacía (`n = 0`): el término inicial es idéntico al final. -/}
  | cero (M : α) : Cadena R 0 M M
  \LCom{/-- Una cadena de `n` pasos seguida de un paso más. -/}
  | paso \{n : Nat\} \{M P N : α\} : Cadena R n M P → R P N → Cadena R (n + 1) M N
\end{Verbatim}
\end{codigolean}
\begin{codigolean}{\ruta{Thesis/Teoria/Lambda.lean}{133}}
\begin{Verbatim}[commandchars=\\\{\}]
\LKw{theorem} concatenar \{r s : Nat\} \{L M N : α\}
    (hLM : Cadena R r L M) (hMN : Cadena R s M N) : Cadena R (r + s) L N := \LKw{by}
  \LCom{-- Inducción sobre la segunda cadena.}
  \LTac{induction} hMN \LKw{with}
  | cero M => \LTac{exact} hLM
  | paso _ hUltimo hipInd => \LTac{exact} Cadena.paso (hipInd hLM) hUltimo
\end{Verbatim}
\end{codigolean}

%---------------------------------------------------------------------
\subsection{Recursión bien fundada: inducción sobre una medida}\label{esq:bienfundada}
\begin{esquema}{Esquema 12. Inducción fuerte sobre la longitud}
\forma «Por inducción fuerte sobre \(|l|\): suponemos el resultado para toda
lista más corta». Se usa cuando las llamadas recursivas no se hacen sobre
subterminos inmediatos (Quicksort llama sobre \emph{filtros} de la cola).
\begin{codigolean}{Plantilla}
\begin{Verbatim}[commandchars=\\\{\}]
\LKw{theorem} t : ∀ l : List α, P l
  | [] => \LKw{by} …
  | p :: cola => \LKw{by}
      \LKw{have} hLt : (f cola).length < (p :: cola).length := …
      \LKw{have} hipInd : P (f cola) := t (f cola)   \LCom{-- llamada recursiva}
      …
  \LKw{termination_by} l => l.length
  \LKw{decreasing_by} \LTac{exact} hLt
\end{Verbatim}
\end{codigolean}
\lectura Un teorema que se invoca a sí mismo es una prueba por inducción. Lean
solo lo acepta si la llamada recursiva es sobre un argumento menor según la
medida declarada en \lean{termination_by}. La obligación correspondiente se
descarga en \lean{decreasing_by}, aquí con los hechos \lean{hLt} ya
probados. Sin esa comprobación el «teorema» podría ser circular.
\donde La definición \lean{quicksort}, \lean{quicksort_permutacion},
\lean{quicksort_ordenada} (medida: longitud); \lean{sustitucion} de
\texttt{SustitucionAlfa} (medida: \lean{tamano}, con \lean{omega}).
\end{esquema}
\begin{codigolean}{\ruta{Thesis/Sort/Quicksort.lean}{159}}
\begin{Verbatim}[commandchars=\\\{\}]
\LKw{def} quicksort : List α → List α
  | [] => []
  | p :: cola =>
      quicksort (cola.filter (\LKw{fun} x => \LTac{decide} (x ≤ p)))
        ++ p :: quicksort (cola.filter (\LKw{fun} x => ! \LTac{decide} (x ≤ p)))
  \LKw{termination_by} l => l.length
  \LKw{decreasing_by}
    \LCom{-- Cada meta compara una partición con p :: cola, no con cola.}
    \LCom{-- Lean añade pruebas de pertenencia con attach; se eliminan de las listas.}
    · \LTac{rw} [List.unattach_filter (l := cola.attach)
          (f := \LKw{fun} x => \LTac{decide} (x.val ≤ p))
          (g := \LKw{fun} x => \LTac{decide} (x ≤ p)) (hf := \LKw{fun} _ _ => \LTac{rfl}),
          List.unattach_attach]
      \LTac{exact} longitud_filtro_lt_cons (\LKw{fun} x => \LTac{decide} (x ≤ p)) p cola
    · \LTac{rw} [List.unattach_filter (l := cola.attach)
          (f := \LKw{fun} x => ! \LTac{decide} (x.val ≤ p))
          (g := \LKw{fun} x => ! \LTac{decide} (x ≤ p)) (hf := \LKw{fun} _ _ => \LTac{rfl}),
          List.unattach_attach]
      \LTac{exact} longitud_filtro_lt_cons (\LKw{fun} x => ! \LTac{decide} (x ≤ p)) p cola
\end{Verbatim}
\end{codigolean}
\begin{codigolean}{\ruta{Thesis/Sort/Quicksort.lean}{144}}
\begin{Verbatim}[commandchars=\\\{\}]
\LKw{theorem} longitud_filtro_lt_cons (q : α → Bool) (p : α) (cola : List α) :
    (cola.filter q).length < (p :: cola).length := \LKw{by}
  \LCom{-- A probar: longitud del filtro < longitud de la lista con pivote.}
  \LCom{-- Método: cota del filtro, seguida de n < n + 1.}
  \LKw{have} hFiltro : (cola.filter q).length ≤ cola.length :=
    longitud_filtro_le q cola
  \LKw{have} hCola : cola.length < (p :: cola).length := \LKw{by}
    \LCom{-- La longitud de p :: cola es cola.length + 1.}
    \LTac{change} cola.length < cola.length + 1
    \LTac{exact} Nat.lt_succ_self cola.length
  \LTac{exact} lt_of_le_of_lt hFiltro hCola
\end{Verbatim}
\end{codigolean}
La prueba de \lean{quicksort_permutacion} se recorre completa en la
sección~\ref{sec:guiada-quicksort}.

%---------------------------------------------------------------------
\subsection{Construcción de derivaciones: una línea, un \texttt{have}}\label{esq:derivaciones}
\begin{esquema}{Esquema 13. Transcribir una derivación en cajas o un árbol de tipado}
\forma Una derivación numerada: cada línea es una fórmula (o un juicio) con su
justificación «regla, líneas citadas».
\begin{codigolean}{Plantilla}
\begin{Verbatim}[commandchars=\\\{\}]
  \LKw{have} l1 : Γ₂ ⊢ₚ φ₁ := Derivable.hip (…pertenencia…)   \LCom{-- premisa/suposición}
  \LKw{have} l2 : Γ₂ ⊢ₚ φ₂ := Derivable.regla l1 …            \LCom{-- regla aplicada}
  \LKw{have} l3 : Γ₁ ⊢ₚ φ₃ := Derivable.implI l2               \LCom{-- cierra caja}
  \LTac{exact} l3
\end{Verbatim}
\end{codigolean}
\lectura Cada \lean{have lN} enuncia el \emph{secuente} completo de la línea
\(N\): la fórmula y también las suposiciones abiertas en ese punto. Abrir una
caja es pasar a \lean{insert φ Γ}; cerrarla es aplicar la regla que descarga
\lean{φ} (\lean{implI}, \lean{negI}, \lean{raa}, \lean{disyE}). Lean comprueba
que la regla cite líneas del contexto correcto.
\donde Capítulo~9: \lean{distributividad_conj_disy}, \lean{contrapositiva_ida},
\lean{contrapositiva_vuelta}, \lean{modus_tollens},
\lean{eliminacion_doble_negacion}, \lean{medio_excluido}. Capítulo~11:
\lean{ej_identidad}, \lean{ej_k}, \lean{ej_eliminacion_doble_negacion},
\lean{introduccion_doble_negacion}, \lean{analisis_por_casos}. Tipado:
\lean{ejemplo_constante}, \lean{ejemplo_aplicacion}, los ejemplos 8.5, 8.13,
8.15 y 8.20 de \texttt{Cubo.lean}.
\end{esquema}
\begin{codigolean}{\ruta{Thesis/Teoria/DeduccionNatural.lean}{204}}
\begin{Verbatim}[commandchars=\\\{\}]
\LKw{theorem} contrapositiva_ida (p q : Formula) :
    (\{p →ₚ q\} : Set Formula) ⊢ₚ ¬ₚq →ₚ ¬ₚp := \LKw{by}
  \LCom{-- Línea 1: p → q — premisa (citada dentro de las dos cajas abiertas).}
  \LKw{have} l1 : insert p (insert (¬ₚq) (\{p →ₚ q\} : Set Formula)) ⊢ₚ p →ₚ q :=
    Derivable.hip (Set.mem_insert_of_mem _ (Set.mem_insert_of_mem _ (Set.mem_singleton
        \LCont{↪ }_)))
  \LCom{-- Línea 2 (abre la caja exterior): ¬q — suposición para →i.}
  \LKw{have} l2 : insert p (insert (¬ₚq) (\{p →ₚ q\} : Set Formula)) ⊢ₚ ¬ₚq :=
    Derivable.hip (Set.mem_insert_of_mem _ (Set.mem_insert _ _))
  \LCom{-- Línea 3 (abre la caja interior): p — suposición para ¬i.}
  \LKw{have} l3 : insert p (insert (¬ₚq) (\{p →ₚ q\} : Set Formula)) ⊢ₚ p :=
    Derivable.hip (Set.mem_insert _ _)
  \LCom{-- Línea 4: q — →e 1, 3.}
  \LKw{have} l4 : insert p (insert (¬ₚq) (\{p →ₚ q\} : Set Formula)) ⊢ₚ q :=
    Derivable.implE l3 l1
  \LCom{-- Línea 5: ⊥ — ¬e 4, 2.}
  \LKw{have} l5 : insert p (insert (¬ₚq) (\{p →ₚ q\} : Set Formula)) ⊢ₚ ⊥ₚ :=
    Derivable.negE l4 l2
  \LCom{-- Línea 6 (cierra la caja interior): ¬p — ¬i 3–5; descarga p.}
  \LKw{have} l6 : insert (¬ₚq) (\{p →ₚ q\} : Set Formula) ⊢ₚ ¬ₚp := Derivable.negI l5
  \LCom{-- Línea 7 (cierra la caja exterior): ¬q → ¬p — →i 2–6; descarga ¬q.}
  \LKw{have} l7 : (\{p →ₚ q\} : Set Formula) ⊢ₚ ¬ₚq →ₚ ¬ₚp := Derivable.implI l6
  \LTac{exact} l7
\end{Verbatim}
\end{codigolean}

La correspondencia con la derivación en cajas de la Proposición~9.9, con
\(\Gamma_0 = \{p \to q\}\), \(\Gamma_1 = \Gamma_0 \cup \{\neg q\}\) y
\(\Gamma_2 = \Gamma_1 \cup \{p\}\):
\begin{center}
\small
\begin{tabular}{@{}cllll@{}}
\toprule
\textbf{Línea} & \textbf{Fórmula} & \textbf{Justificación} & \textbf{Término Lean} & \textbf{Contexto} \\
\midrule
1 & \(p \to q\) & premisa & \lean{hip (…of_mem (…of_mem (mem_singleton)))} & \(\Gamma_2\) \\
2 & \(\quad\neg q\) & suposición (\(\to\)i) & \lean{hip (mem_insert_of_mem _ (mem_insert _ _))} & \(\Gamma_2\) \\
3 & \(\qquad p\) & suposición (\(\neg\)i) & \lean{hip (mem_insert _ _)} & \(\Gamma_2\) \\
4 & \(\qquad q\) & \(\to\)e 1, 3 & \lean{implE l3 l1} & \(\Gamma_2\) \\
5 & \(\qquad \bot\) & \(\neg\)e 4, 2 & \lean{negE l4 l2} & \(\Gamma_2\) \\
6 & \(\quad\neg p\) & \(\neg\)i 3--5 & \lean{negI l5} & \(\Gamma_1\) \\
7 & \(\neg q \to \neg p\) & \(\to\)i 2--6 & \lean{implI l6} & \(\Gamma_0\) \\
\bottomrule
\end{tabular}
\end{center}

\textbf{Cómo leer las pruebas de pertenencia.} El contexto
\lean{insert p (insert (¬ₚq) {p →ₚ q})} es una pila de tres capas.
\lean{Set.mem_insert _ _} toma la capa superior; cada
\lean{Set.mem_insert_of_mem _} baja una capa; \lean{Set.mem_singleton _} toma
la base. La línea 1 baja dos capas y toma la base: es la premisa original,
citada desde dentro de dos cajas. Estas pruebas son contabilidad: responden a
«¿esta fórmula está disponible aquí?».

En los cálculos de tipos el esquema es el mismo; cada paso numerado del texto
es un \lean{have pasoN} o, en \texttt{Cubo.lean}, un teorema propio:
\begin{codigolean}{\ruta{Thesis/Teoria/TiposSimples.lean}{125}}
\begin{Verbatim}[commandchars=\\\{\}]
\LKw{theorem} ejemplo_constante (σ τ : Tipo) :
    [] ⊢ .abs "x" σ (.abs "y" τ (.var "x")) : σ ⟶ (τ ⟶ σ) := \LKw{by}
  \LCom{-- Contexto inicial x:σ, y:τ; sus sujetos son distintos.}
  \LKw{have} hContexto : BienFormado [("x", σ), ("y", τ)] := \LKw{by}
    \LTac{simp} [BienFormado]
  \LCom{-- Paso (1), regla (variable): x:σ, y:τ ⊢ x : σ, pues (x:σ) está en el contexto.}
  \LKw{have} paso1 : [("x", σ), ("y", τ)] ⊢ .var "x" : σ :=
    Deriva.var hContexto (List.mem_cons_self)
  \LCom{-- Paso (2), regla (abstracción) sobre (1): x:σ ⊢ λy:τ. x : τ → σ.}
  \LKw{have} paso2 : [("x", σ)] ⊢ .abs "y" τ (.var "x") : τ ⟶ σ :=
    Deriva.abst (Γ := [("x", σ)]) paso1
  \LCom{-- Paso (3), regla (abstracción) sobre (2): ⊢ λx:σ. λy:τ. x : σ → (τ → σ).}
  \LTac{exact} Deriva.abst (Γ := []) paso2
\end{Verbatim}
\end{codigolean}
\begin{codigolean}{\ruta{Thesis/Teoria/Cubo.lean}{382}}
\begin{Verbatim}[commandchars=\\\{\}]
\LKw{theorem} completa_paso2 : Deriva (ctxBeta ,, "α" ∶ ★) α ★ := \LKw{by}
  \LCom{-- Condición lateral: α ∉ (β : ∗), es decir, α y β son nombres distintos.}
  \LKw{have} hFresca : "α" ∉ ctxBeta.dominio := \LKw{by} \LTac{simp} [Contexto.dominio]
  \LCom{-- Paso 2: (variable) de 1.}
  \LTac{exact} Deriva.variable completa_paso1 hFresca
\end{Verbatim}
\end{codigolean}
\begin{codigolean}{\ruta{Thesis/Teoria/Cubo.lean}{412}}
\begin{Verbatim}[commandchars=\\\{\}]
\LKw{theorem} derivacion_completa : Deriva ctxBeta (.apl constructorFlecha β) ★ :=
  \LCom{-- Paso 7: (aplicación) de 5, 6.}
  Deriva.aplicacion completa_paso5 completa_paso6
\end{Verbatim}
\end{codigolean}
Las condiciones laterales (\(\alpha \notin \Gamma\), sujetos distintos) se
cierran con \lean{simp [Contexto.dominio]} o \lean{simp [BienFormado]}: son
comprobaciones decidibles sobre listas finitas de nombres, no razonamiento.

\begin{nota}[El orden de las premisas]
En \lean{Derivable} (capítulo~9) la regla \lean{implE} recibe primero
\(\varphi\) y luego \(\varphi \to \psi\), como en la tesis («\(\to\)e 3, 1»
en el orden de Huth y Ryan). En \lean{ND} (capítulo~11), \lean{impE} recibe
primero la implicación. El archivo \texttt{DeduccionNatural.lean} transcribe
la Proposición~9.9 en ambos cálculos (\lean{contrapositiva_ida_ND}).
\end{nota}

%---------------------------------------------------------------------
\subsection{Composición de lemas}\label{esq:composicion}
\begin{esquema}{Esquema 14. El teorema principal como ensamblaje}
\forma Un resultado grande se obtiene encadenando lemas: «por el Lema A,
\ldots; por el Lema B aplicado a eso, \ldots».
\begin{codigolean}{Plantilla}
\begin{Verbatim}[commandchars=\\\{\}]
  \LKw{have} h₁ : … := lemaA …
  \LKw{have} h₂ : … := lemaB h₁ …
  \LTac{exact} lemaC h₂
\end{Verbatim}
\end{codigolean}
\lectura La prueba del teorema principal suele ser corta: el trabajo está en
los lemas. Al exponer, conviene presentar primero el ensamblaje y después cada
pieza.
\donde \lean{completitud_ND} (Kalmár + descarga),
\lean{correccion_completitud}, \lean{unicidad_forma_normal}
(Lema~7.32 + Corolario~7.43 + forma normal), \lean{quicksort_correcto},
\lean{lema_7_6}, \lean{lema_7_31}, \lean{lema_7_32}.
\end{esquema}
\begin{codigolean}{\ruta{Thesis/Prop/Completeness.lean}{256}}
\begin{Verbatim}[commandchars=\\\{\}]
\LKw{theorem} completitud_ND (φ : Formula) (h : EsTautologia φ) : ND (∅ : Set Formula) φ := \LKw{by}
  \LCom{-- A probar: ND ∅ φ. Método: Kalmár para cada valuación, seguido de descarga.}
  \LKw{let} listaAtomos := (atomos φ).dedup
  \LKw{have} hSinRepeticiones : listaAtomos.Nodup := List.nodup_dedup _
  \LKw{have} hCobertura : atomos φ ⊆ listaAtomos := \LKw{by}
    \LTac{intro} a ha
    \LTac{exact} List.mem_dedup.mpr ha
  \LKw{have} hUniversal : ∀ v, ND (contextoLiterales v listaAtomos) φ := \LKw{by}
    \LTac{intro} v
    \LKw{have} hVerdadero : evaluar v φ := h v
    \LKw{have} hPositiva : evaluar v φ → ND (contextoLiterales v listaAtomos) φ :=
      (kalmar v listaAtomos φ hCobertura).left
    \LTac{exact} hPositiva hVerdadero
  \LTac{exact} descarga listaAtomos hSinRepeticiones hUniversal
\end{Verbatim}
\end{codigolean}
\begin{codigolean}{\ruta{Thesis/Teoria/Lambda.lean}{400}}
\begin{Verbatim}[commandchars=\\\{\}]
\LKw{theorem} unicidad_forma_normal (hCR : Confluente sust) \{M N₁ N₂ : Termino\}
    (h₁ : EsFormaNormalDe sust N₁ M) (h₂ : EsFormaNormalDe sust N₂ M) : N₁ = N₂ := \LKw{by}
  \LCom{-- `M =β N₁` y `M =β N₂` con `N₁`, `N₂` en forma normal β.}
  \LTac{obtain} ⟨hNormal₁, hMN₁⟩ := h₁
  \LTac{obtain} ⟨hNormal₂, hMN₂⟩ := h₂
  \LCom{-- Por simetría y transitividad (Lema 7.32), `N₁ =β N₂`.}
  \LKw{have} hN₁N₂ : IgualBeta sust N₁ N₂ :=
    igualBeta_transitiva (igualBeta_simetrica hMN₁) hMN₂
  \LCom{-- Por el Corolario 7.43 existe `L` con `N₁ ↠β L` y `N₂ ↠β L`.}
  \LTac{obtain} ⟨L, hN₁L, hN₂L⟩ := church_rosser_conversion hCR hN₁N₂
  \LCom{-- Una forma normal solo se reduce a sí misma: `N₁ ≡ L` y `N₂ ≡ L`.}
  \LKw{have} hIgual₁ : N₁ = L := formaNormal_reduce_igual hNormal₁ hN₁L
  \LKw{have} hIgual₂ : N₂ = L := formaNormal_reduce_igual hNormal₂ hN₂L
  \LCom{-- De `N₁ ≡ L ≡ N₂` se concluye `N₁ ≡ N₂`.}
  \LTac{rw} [hIgual₁, hIgual₂]
\end{Verbatim}
\end{codigolean}

%---------------------------------------------------------------------
\subsection{Resultados citados como hipótesis explícitas}\label{esq:hipotesis}
\begin{esquema}{Esquema 15. Hacer visible lo que no se demuestra}
\forma «Suponiendo el Teorema~X (que no demostramos), se cumple Y».
\begin{codigolean}{Plantilla}
\begin{Verbatim}[commandchars=\\\{\}]
\LKw{def} EnunciadoX : Prop := …
\LKw{theorem} y (hX : EnunciadoX) … : … := \LKw{by} …     \LCom{-- hX es una premisa}
\end{Verbatim}
\end{codigolean}
\lectura En vez de declarar un \lean{axiom}, el resultado no demostrado se
define como una proposición y se pide como argumento. El teorema final queda
condicionado a ella de forma visible, y la auditoría de axiomas no cambia.
\donde El Teorema~7.37 (Church--Rosser), que la tesis cita sin demostración,
es \lean{Confluente} y entra como \lean{hCR} en \lean{church_rosser_conversion}
y \lean{unicidad_forma_normal}. La sustitución de la reducción~\(\beta\) es un
parámetro (\lean{variable (sust : …)}), porque esos resultados no usan sus
propiedades. La convención de variables de Barendregt es la hipótesis
\lean{ConvencionVariables} en \texttt{Sustitucion.lean}.
\end{esquema}
\begin{codigolean}{\ruta{Thesis/Teoria/Lambda.lean}{299}}
\begin{Verbatim}[commandchars=\\\{\}]
\LKw{def} Confluente : Prop :=
  ∀ M N₁ N₂ : Termino, ReduceBeta sust M N₁ → ReduceBeta sust M N₂ →
    ∃ L : Termino, ReduceBeta sust N₁ L ∧ ReduceBeta sust N₂ L
\end{Verbatim}
\end{codigolean}
\begin{codigolean}{\ruta{Thesis/Teoria/Lambda.lean}{377}}
\begin{Verbatim}[commandchars=\\\{\}]
\LKw{theorem} church_rosser_conversion (hCR : Confluente sust) \{M N : Termino\}
    (h : IgualBeta sust M N) :
    ∃ L : Termino, ReduceBeta sust M L ∧ ReduceBeta sust N L := \LKw{by}
  \LTac{obtain} ⟨k, hCadena⟩ := h
  \LCom{-- Inducción sobre la longitud `k` de la cadena `M ≡ M₀ ↔β M₁ ↔β ⋯ ↔β M_k ≡ N`.}
  \LTac{induction} hCadena \LKw{with}
  | cero M =>
      \LCom{-- `k = 0`: basta `L ≡ M`.}
      \LTac{exact} ⟨M, reduceBeta_reflexiva M, reduceBeta_reflexiva M⟩
  | @paso k M Mk1 Mk _ hUltimo hipInd =>
      \LCom{-- La hipótesis de inducción da `L'` con `M ↠β L'` y `M_\{k-1\} ↠β L'`.}
      \LTac{obtain} ⟨L', hML', hMk1L'⟩ := hipInd
      \LTac{rcases} hUltimo \LKw{with} hDirecto | hInverso
      · \LCom{-- `M_\{k-1\} →β M_k`: el Teorema 7.37 aplicado a `M_\{k-1\}`, que se reduce a `L'`}
        \LCom{-- y a `M_k`, da `L` con `L' ↠β L` y `M_k ↠β L`.}
        \LTac{obtain} ⟨L, hL'L, hMkL⟩ := hCR Mk1 L' Mk hMk1L' (reduceBeta_de_paso hDirecto)
        \LCom{-- Componiendo, `M ↠β L' ↠β L`.}
        \LTac{exact} ⟨L, reduceBeta_transitiva hML' hL'L, hMkL⟩
      · \LCom{-- `M_k →β M_\{k-1\}`: entonces `M_k ↠β L'` y se toma `L ≡ L'`.}
        \LTac{exact} ⟨L', hML', reduceBeta_transitiva (reduceBeta_de_paso hInverso) hMk1L'⟩
\end{Verbatim}
\end{codigolean}
Las dos ramas del \lean{rcases} son los dos sentidos posibles del último paso
de conversión:
\begin{estadolean}{Estado antes}
\begin{Verbatim}[commandchars=\\\{\}]
\LCom{-- (… 1 hipótesis omitidas …)}
hCR : Confluente sust
\LCom{-- (… 5 hipótesis omitidas …)}
hUltimo : PasoConversion sust Mk1 Mk
\LCom{-- (… 1 hipótesis omitidas …)}
hML' : ReduceBeta sust M L'
hMk1L' : ReduceBeta sust Mk1 L'
⊢ ∃ L, ReduceBeta sust M L ∧ ReduceBeta sust Mk L
\end{Verbatim}
\end{estadolean}
\tacticaentre{rcases hUltimo with hDirecto | hInverso}
\begin{estadolean}{Estado después}
\begin{Verbatim}[commandchars=\\\{\}]
\LCom{-- meta 1: case paso.inl}
\LCom{-- (… 1 hipótesis omitidas …)}
hCR : Confluente sust
\LCom{-- (… 6 hipótesis omitidas …)}
hML' : ReduceBeta sust M L'
hMk1L' : ReduceBeta sust Mk1 L'
hDirecto : PasoBeta sust Mk1 Mk
⊢ ∃ L, ReduceBeta sust M L ∧ ReduceBeta sust Mk L

\LCom{-- meta 2: case paso.inr}
\LCom{-- (… 1 hipótesis omitidas …)}
hCR : Confluente sust
\LCom{-- (… 6 hipótesis omitidas …)}
hML' : ReduceBeta sust M L'
hMk1L' : ReduceBeta sust Mk1 L'
hInverso : PasoBeta sust Mk Mk1
⊢ ∃ L, ReduceBeta sust M L ∧ ReduceBeta sust Mk L
\end{Verbatim}
\end{estadolean}

%---------------------------------------------------------------------
\subsection{Transporte a través de una representación}\label{esq:transporte}
\begin{esquema}{Esquema 16. Probar en una representación más cómoda y traducir}
\forma «Traducimos al lenguaje B, donde el resultado es fácil, y la
traducción preserva lo necesario».
\begin{codigolean}{Plantilla}
\begin{Verbatim}[commandchars=\\\{\}]
  \LTac{calc} tr (lado izquierdo)
      = operación_B (tr …) := lema_de_traducción
    _ = operación_B' (tr …) := resultado_en_B
    _ = tr (lado derecho) := lema_de_traducción
\end{Verbatim}
\end{codigolean}
\donde \lean{lema_sustitucion_alfa} (Lema~7.23 módulo~\(\alpha\)): la
\(\alpha\)-equivalencia se define como igualdad de traducciones a índices de
De~Bruijn, donde la sustitución no captura. \lean{traduccion_derivable_de_tautologia}
(capítulo~11): la completitud se valida de forma independiente con la
biblioteca Foundation, traduciendo fórmulas y usando \lean{evaluar_traduccion}.
\end{esquema}
\begin{codigolean}{\ruta{Thesis/Teoria/SustitucionAlfa.lean}{515}}
\begin{Verbatim}[commandchars=\\\{\}]
\LKw{def} AlfaEq (M N : Termino) : Prop :=
  aDeBruijn [] M = aDeBruijn [] N
\end{Verbatim}
\end{codigolean}
\begin{codigolean}{\ruta{Thesis/Teoria/SustitucionAlfa.lean}{769}}
\begin{Verbatim}[commandchars=\\\{\}]
\LKw{theorem} lema_sustitucion_alfa (M N L : Termino) \{x y : String\}
    (hxy : x ≠ y) (hxL : x ∉ libres L) :
    M⟦x := N⟧⟦y := L⟧ ≡α M⟦y := L⟧⟦x := N⟦y := L⟧⟧ := \LKw{by}
  \LCom{-- `x ∉ FV(L)`: en la traducción de `L` no hay ninguna `x` que sustituir.}
  \LKw{have} hxl : ∀ u, sustLibre (aDeBruijn [] L) x u = aDeBruijn [] L := \LKw{fun} u =>
    sustLibre_aDeBruijn_sin_libre u L [] (\LKw{fun} h => absurd h hxL)
  \LCom{-- Se traducen las cuatro sustituciones y se aplica el lema para `TerminoDB`.}
  \LTac{calc} aDeBruijn [] (M⟦x := N⟧⟦y := L⟧)
      = sustLibre (sustLibre (aDeBruijn [] M) x (aDeBruijn [] N)) y (aDeBruijn [] L) :=
          \LCont{↪ }\LKw{by}
        \LTac{rw} [aDeBruijn_sustitucion, aDeBruijn_sustitucion]
    _ = sustLibre (sustLibre (aDeBruijn [] M) y (aDeBruijn [] L)) x
          (sustLibre (aDeBruijn [] N) y (aDeBruijn [] L)) :=
        sustLibre_conmuta hxy _ _ hxl _
    _ = aDeBruijn [] (M⟦y := L⟧⟦x := N⟦y := L⟧⟧) := \LKw{by}
        \LTac{rw} [aDeBruijn_sustitucion, aDeBruijn_sustitucion, aDeBruijn_sustitucion]
\end{Verbatim}
\end{codigolean}
El \lean{calc} tiene tres pasos: traducir el lado izquierdo, aplicar el lema
para la representación sin nombres (\lean{sustLibre_conmuta}, por inducción
estructural con los tres subcasos de la tesis) y traducir de vuelta.

%---------------------------------------------------------------------
\subsection{Validar la definición: equivalencia con la del texto}\label{esq:validar}
\begin{esquema}{Esquema 17. La formalización dice lo que dice la tesis}
\forma Cuando la definición formal difiere de la del texto (por comodidad),
se demuestra que ambas son equivalentes.
\begin{codigolean}{Plantilla}
\begin{Verbatim}[commandchars=\\\{\}]
\LKw{theorem} defFormal_iff_defTexto : DefFormal x ↔ DefTexto x :=
  ⟨formal_implica_texto, texto_implica_formal⟩
\end{Verbatim}
\end{codigolean}
\lectura Lean garantiza que la prueba es correcta \emph{para las definiciones
escritas}. Que esas definiciones capturen las de la tesis es una afirmación
que hay que justificar, y la forma más sólida es un teorema de equivalencia.
\donde \lean{alfaEq_iff_alfaTesis} (\(\equiv_\alpha\) vía De~Bruijn coincide
con la relación generada por las reglas de la Definición~7.16);
\lean{formaNormal_sin_paso} y \lean{paso_de_redex} (las dos lecturas de
«forma normal» de la Definición~7.33 coinciden); \lean{Deriva.bienFormado}
(todo contexto derivable tiene sujetos distintos, como pide la
Definición~7.51); los tres \lean{example} de \texttt{Cubo.lean} que comprueban
que las derivaciones abreviada y completa del Ejemplo~8.20 enuncian los mismos
juicios.
\end{esquema}
\begin{codigolean}{\ruta{Thesis/Teoria/SustitucionAlfa.lean}{598}}
\begin{Verbatim}[commandchars=\\\{\}]
\LKw{inductive} AlfaTesis : Termino → Termino → Prop
  | renombre (x y : String) (M : Termino) :
      y ∉ libres M → y ∉ ligadas M → AlfaTesis (lam x M) (lam y (renombrar M x y))
  | appIzq \{M N : Termino\} (L : Termino) : AlfaTesis M N → AlfaTesis (app M L) (app N L)
  | appDer \{M N : Termino\} (L : Termino) : AlfaTesis M N → AlfaTesis (app L M) (app L N)
  | abs \{M N : Termino\} (z : String) : AlfaTesis M N → AlfaTesis (lam z M) (lam z N)
  | refl (M : Termino) : AlfaTesis M M
  | simm \{M N : Termino\} : AlfaTesis M N → AlfaTesis N M
  | trans \{L M N : Termino\} : AlfaTesis L M → AlfaTesis M N → AlfaTesis L N
\end{Verbatim}
\end{codigolean}
\begin{codigolean}{\ruta{Thesis/Teoria/SustitucionAlfa.lean}{757}}
\begin{Verbatim}[commandchars=\\\{\}]
\LKw{theorem} alfaEq_iff_alfaTesis \{M N : Termino\} : M ≡α N ↔ AlfaTesis M N :=
  ⟨alfaEq_implica_alfaTesis, alfaTesis_implica_alfaEq⟩
\end{Verbatim}
\end{codigolean}

%=====================================================================
\section{Lecturas guiadas con estados reales}\label{sec:guiadas}
%=====================================================================

En esta sección se siguen tres pruebas táctica por táctica. Cada recuadro
amarillo es el estado que Lean calcula en ese punto.

\subsection{Debilitamiento, caso \texttt{impI}}\label{sec:guiada-debilitamiento}

Enunciado: si \(\Gamma \vdash \varphi\) y \(\Gamma \subseteq \Delta\), entonces
\(\Delta \vdash \varphi\). Tras \lean{induction d with}, la rama \lean{impI}
empieza así. La derivación terminaba con \(\to\)I, de modo que
\(\varphi\) es ahora \(\varphi \to \psi\) y la hipótesis inductiva habla de la
premisa \(\Gamma, \varphi \vdash \psi\):
\begin{estadolean}{Estado antes de \texttt{intro Δ hsub}}
\begin{Verbatim}[commandchars=\\\{\}]
Γ : Set Formula
φ : Formula
ψ : Formula
hipInd : ∀ \{Δ : Set Formula\}, insert φ Γ ⊆ Δ → ND Δ ψ
⊢ ∀ \{Δ : Set Formula\}, Γ ⊆ Δ → ND Δ (φ ⟶ ψ)
\end{Verbatim}
\end{estadolean}
Se fija un \(\Delta\) mayor:
\begin{estadolean}{Estado después de \texttt{intro Δ hsub}}
\begin{Verbatim}[commandchars=\\\{\}]
Γ : Set Formula
φ : Formula
ψ : Formula
hipInd : ∀ \{Δ : Set Formula\}, insert φ Γ ⊆ Δ → ND Δ ψ
Δ : Set Formula
hsub : Γ ⊆ Δ
⊢ ND Δ (φ ⟶ ψ)
\end{Verbatim}
\end{estadolean}
Se amplía la inclusión con \(\varphi\) en ambos lados y se aplica la hipótesis
inductiva a ese contexto ampliado:
\begin{estadolean}{Estado después de \texttt{have dCuerpo : ND (insert φ Δ) ψ := hipInd hAmpliado}}
\begin{Verbatim}[commandchars=\\\{\}]
Γ : Set Formula
φ : Formula
ψ : Formula
hipInd : ∀ \{Δ : Set Formula\}, insert φ Γ ⊆ Δ → ND Δ ψ
Δ : Set Formula
hsub : Γ ⊆ Δ
hAmpliado : insert φ Γ ⊆ insert φ Δ
dCuerpo : ND (insert φ Δ) ψ
⊢ ND Δ (φ ⟶ ψ)
\end{Verbatim}
\end{estadolean}
Por último, \lean{ND.impI dCuerpo} descarga \(\varphi\) y cierra la meta. En
papel: «Si la última regla es \(\to\)I, la premisa es \(\Gamma, \varphi \vdash
\psi\). Como \(\Gamma \cup \{\varphi\} \subseteq \Delta \cup \{\varphi\}\), por
hipótesis inductiva \(\Delta, \varphi \vdash \psi\), y por \(\to\)I,
\(\Delta \vdash \varphi \to \psi\)».

\subsection{Lema de Kalmár, caso atómico}

El enunciado completo es una conjunción de dos implicaciones, bajo la hipótesis
de que \lean{listaAtomos} cubre los átomos de \(\varphi\):
\begin{estadolean}{Estado antes de \texttt{intro φ}}
\begin{Verbatim}[commandchars=\\\{\}]
v : Valuacion
listaAtomos : List String
⊢ ∀ (φ : Formula),
  atomos φ ⊆ listaAtomos →
    (evaluar v φ → ND (contextoLiterales v listaAtomos) φ) ∧ (¬evaluar v φ → ND
        \LCont{↪ }(contextoLiterales v listaAtomos) (~φ))
\end{Verbatim}
\end{estadolean}
En el caso \lean{atom s}, tras obtener \lean{hs : s ∈ listaAtomos} se divide
la conjunción:
\begin{estadolean}{Estado antes}
\begin{Verbatim}[commandchars=\\\{\}]
\LCom{-- (… 3 hipótesis omitidas …)}
hsub : atomos (Formula.atom s) ⊆ listaAtomos
\LCom{-- (… 1 hipótesis omitidas …)}
hs : s ∈ listaAtomos
⊢ (evaluar v (Formula.atom s) → ND (contextoLiterales v listaAtomos) (Formula.atom s))
    \LCont{↪ }∧
  (¬evaluar v (Formula.atom s) → ND (contextoLiterales v listaAtomos) (~Formula.atom
      \LCont{↪ }s))
\end{Verbatim}
\end{estadolean}
\tacticaentre{refine ⟨?\_, ?\_⟩}
\begin{estadolean}{Estado después}
\begin{Verbatim}[commandchars=\\\{\}]
\LCom{-- meta 1: case atom.refine_1}
\LCom{-- (… 3 hipótesis omitidas …)}
hsub : atomos (Formula.atom s) ⊆ listaAtomos
\LCom{-- (… 1 hipótesis omitidas …)}
hs : s ∈ listaAtomos
⊢ evaluar v (Formula.atom s) → ND (contextoLiterales v listaAtomos) (Formula.atom s)

\LCom{-- meta 2: case atom.refine_2}
\LCom{-- (… 3 hipótesis omitidas …)}
hsub : atomos (Formula.atom s) ⊆ listaAtomos
\LCom{-- (… 1 hipótesis omitidas …)}
hs : s ∈ listaAtomos
⊢ ¬evaluar v (Formula.atom s) → ND (contextoLiterales v listaAtomos) (~Formula.atom s)
\end{Verbatim}
\end{estadolean}
En la parte positiva, si \(v(s)\) es verdadero, el literal de \(s\) es el
propio átomo \(s\); por \lean{pertenencia_literal} está en el contexto de
literales, y la regla \lean{ND.hyp} lo deriva. La parte negativa es simétrica
con \(\neg s\).

\subsection{Quicksort produce una permutación}\label{sec:guiada-quicksort}
\begin{codigolean}{\ruta{Thesis/Sort/Quicksort.lean}{195}}
\begin{Verbatim}[commandchars=\\\{\}]
\LKw{theorem} quicksort_permutacion : ∀ l : List α, quicksort l ~ l
  | [] => \LKw{by}
      \LCom{-- A probar: quicksort [] ~ []. Reescritura y reflexividad.}
      \LTac{simpa} \LTac{only} [quicksort_vacia] \LTac{using} (List.Perm.refl ([] : List α))
  | p :: cola => \LKw{by}
      \LCom{-- A probar: quicksort (p :: cola) ~ p :: cola.}
      \LCom{-- Método: inducción por longitud; las particiones son menores.}
      \LKw{let} menoresIguales := cola.filter (\LKw{fun} x => \LTac{decide} (x ≤ p))
      \LKw{let} mayores := cola.filter (\LKw{fun} x => ! \LTac{decide} (x ≤ p))
      \LKw{have} hMenoresIgualesLt : menoresIguales.length < (p :: cola).length :=
        longitud_filtro_lt_cons (\LKw{fun} x => \LTac{decide} (x ≤ p)) p cola
      \LKw{have} hMayoresLt : mayores.length < (p :: cola).length :=
        longitud_filtro_lt_cons (\LKw{fun} x => ! \LTac{decide} (x ≤ p)) p cola
      \LCom{-- Las cotas anteriores legitiman estas dos instancias inductivas.}
      \LKw{have} hipIndMenoresIguales : quicksort menoresIguales ~ menoresIguales :=
          \LCont{↪ }quicksort_permutacion menoresIguales
      \LKw{have} hipIndMayores : quicksort mayores ~ mayores := quicksort_permutacion mayores
      \LKw{have} hParticion : menoresIguales ++ mayores ~ cola :=
        particion_filtros_permutacion (\LKw{fun} x => \LTac{decide} (x ≤ p)) cola
      \LCom{-- 1. Añadir el mismo pivote conserva la permutación derecha.}
      \LKw{have} hDerecha : p :: quicksort mayores ~ p :: mayores := hipIndMayores.cons p
      \LCom{-- 2. Concatenar las dos permutaciones conserva ambas listas.}
      \LKw{have} hConcatenacion : quicksort menoresIguales ++ (p :: quicksort mayores)
          ~ menoresIguales ++ (p :: mayores) := hipIndMenoresIguales.append hDerecha
      \LCom{-- 3. perm_middle desplaza el pivote al principio.}
      \LKw{have} hPivote : menoresIguales ++ (p :: mayores) ~ p :: (menoresIguales ++ mayores)
          \LCont{↪ }:=
        List.perm_middle
      \LCom{-- 4. La congruencia de :: se aplica a la partición anterior.}
      \LKw{have} hCons : p :: (menoresIguales ++ mayores) ~ p :: cola := hParticion.cons p
      \LCom{-- Reescritura de quicksort, seguida de las tres permutaciones.}
      \LTac{rw} [quicksort_cons]
      \LTac{change} quicksort menoresIguales ++ (p :: quicksort mayores) ~ p :: cola
      \LTac{calc}
        quicksort menoresIguales ++ (p :: quicksort mayores)
            ~ menoresIguales ++ (p :: mayores) := hConcatenacion
        _ ~ p :: (menoresIguales ++ mayores) := hPivote
        _ ~ p :: cola := hCons
  \LKw{termination_by} l => l.length
  \LKw{decreasing_by}
    · \LTac{exact} hMenoresIgualesLt
    · \LTac{exact} hMayoresLt
\end{Verbatim}
\end{codigolean}
Antes de usar la ecuación de \lean{quicksort} la meta todavía menciona
\lean{quicksort (p :: cola)}; tras \lean{rw [quicksort_cons]} y \lean{change}
queda en la forma que admite el \lean{calc}:
\begin{estadolean}{Estado antes de \texttt{calc}}
\begin{Verbatim}[commandchars=\\\{\}]
α : Type u_1
p : α
cola : List α
menoresIguales : List α := filter (\LKw{fun} x => \LTac{decide} (x ≤ p)) cola
mayores : List α := filter (\LKw{fun} x => !decide (x ≤ p)) cola
hMenoresIgualesLt : menoresIguales.length < (p :: cola).length
hMayoresLt : mayores.length < (p :: cola).length
hipIndMenoresIguales : quicksort menoresIguales ~ menoresIguales
hipIndMayores : quicksort mayores ~ mayores
hParticion : menoresIguales ++ mayores ~ cola
hDerecha : p :: quicksort mayores ~ p :: mayores
hConcatenacion : quicksort menoresIguales ++ p :: quicksort mayores ~ menoresIguales
    \LCont{↪ }++ p :: mayores
hPivote : menoresIguales ++ p :: mayores ~ p :: (menoresIguales ++ mayores)
hCons : p :: (menoresIguales ++ mayores) ~ p :: cola
⊢ quicksort menoresIguales ++ p :: quicksort mayores ~ p :: cola
\end{Verbatim}
\end{estadolean}
El \lean{calc} es la cadena de permutaciones del capítulo~12:
\[
\mathit{qs}(L) \mathbin{+\!+} (p :: \mathit{qs}(G))
\;\sim\; L \mathbin{+\!+} (p :: G)
\;\sim\; p :: (L \mathbin{+\!+} G)
\;\sim\; p :: \mathit{cola},
\]
con \(L\) = \lean{menoresIguales} y \(G\) = \lean{mayores}. El primer paso usa
las dos hipótesis inductivas, el segundo mueve el pivote al frente
(\lean{List.perm_middle}) y el tercero usa que los dos filtros reparten la cola
(\lean{particion_filtros_permutacion}). Las dos obligaciones de
\lean{decreasing_by} son exactamente \lean{hMenoresIgualesLt} y
\lean{hMayoresLt}: lo que autoriza llamar a \lean{quicksort_permutacion} sobre
los filtros.

%=====================================================================
\section{Mapa de la tesis}\label{sec:mapa}
%=====================================================================

Números de esquema: 1 directa, 2 hacia adelante, 3 bicondicional,
4 conjunción, 5 casos de hipótesis, 6 \lean{by_cases}, 7 negación/absurdo,
8 existencial, 9 ecuacional, 10 inducción estructural, 11 inducción sobre
derivaciones, 12 recursión bien fundada, 13 construcción de derivaciones,
14 composición, 15 hipótesis explícita, 16 transporte, 17 validación.

{\small
\renewcommand{\arraystretch}{1.18}
\begin{longtable}{@{}>{\raggedright\arraybackslash}p{0.19\textwidth}>{\raggedright\arraybackslash}p{0.33\textwidth}>{\raggedright\arraybackslash}p{0.11\textwidth}>{\raggedright\arraybackslash}p{0.29\textwidth}@{}}
\toprule
\textbf{Resultado} & \textbf{Teorema Lean} & \textbf{Esquemas} & \textbf{Archivo} \\
\midrule
\endhead
\multicolumn{4}{@{}l}{\textit{Capítulo 7: cálculo \(\lambda\) sin tipos}} \\
Lema 7.6 & \lean{sub_reflexiva}, \lean{sub_transitiva}, \lean{lema_7_6} & 5, 10, 4 & \texttt{Teoria/Lambda} \\
Lema 7.31 & \lean{reduceBeta_de_paso}, \lean{reduceBeta_reflexiva}, \lean{reduceBeta_transitiva} & 8, 11, 4 & \texttt{Teoria/Lambda} \\
Lema 7.32 & \lean{igualBeta_de_reduceBeta}, \lean{igualBeta_simetrica}, \lean{igualBeta_transitiva} & 5, 8, 11 & \texttt{Teoria/Lambda} \\
Corolario 7.43 & \lean{church_rosser_conversion} & 11, 5, 8, 15 & \texttt{Teoria/Lambda} \\
Corolario 7.44 & \lean{unicidad_forma_normal} & 14, 15, 9 & \texttt{Teoria/Lambda} \\
Definición 7.33 & \lean{formaNormal_sin_paso}, \lean{paso_de_redex} & 11, 10, 17 & \texttt{Teoria/Lambda} \\
Lema 7.23 (literal) & \lean{lema_sustitucion} & 10, 6, 9, 15 & \texttt{Teoria/Sustitucion} \\
Lema 7.23 (mód.~\(\alpha\)) & \lean{lema_sustitucion_alfa}, \lean{lema_sustitucion_tesis} & 16, 9, 10, 17 & \texttt{Teoria/SustitucionAlfa} \\
Definición 7.16 & \lean{alfaEq_iff_alfaTesis} & 17, 11 & \texttt{Teoria/SustitucionAlfa} \\
\multicolumn{4}{@{}l}{\textit{Capítulo 7: tipos simples}} \\
Ejemplos 7.52, 7.53 & \lean{ejemplo_constante}, \lean{ejemplo_aplicacion} & 13 & \texttt{Teoria/TiposSimples} \\
Lema 7.57 & \lean{variables_libres} & 11, 9, 5, 7 & \texttt{Teoria/TiposSimples} \\
Definición 7.51 & \lean{Deriva.bienFormado} & 11, 17 & \texttt{Teoria/TiposSimples} \\
\multicolumn{4}{@{}l}{\textit{Capítulo 8: cubo \(\lambda\)}} \\
Ejemplo 8.5 & \lean{identidad_polimorfica} & 13 & \texttt{Teoria/Cubo} \\
Ejemplos 8.13, 8.15 & \lean{variable_termino}, \lean{debilitamiento_a}\ldots\lean{_d} & 13 & \texttt{Teoria/Cubo} \\
Ejemplo 8.20 & \lean{derivacion_completa}, \lean{derivacion_abreviada} & 13, 17 & \texttt{Teoria/Cubo} \\
Regla (conversión) & \lean{conversion_ejemplo}, \lean{constructor_beta} & 13, 9 & \texttt{Teoria/Cubo} \\
\multicolumn{4}{@{}l}{\textit{Capítulo 9: deducción natural}} \\
Proposición 9.6 & \lean{distributividad_conj_disy} & 13 & \texttt{Teoria/DeduccionNatural} \\
Proposición 9.9 & \lean{contrapositiva} (y versiones \lean{_ND}) & 13, 4 & \texttt{Teoria/DeduccionNatural} \\
MT, LEM, \(\neg\neg\)e & \lean{modus_tollens}, \lean{medio_excluido}, \lean{eliminacion_doble_negacion} & 13 & \texttt{Teoria/DeduccionNatural} \\
Debilitamiento & \lean{debilitamiento}, \lean{reiterar} & 11 & \texttt{Teoria/DeduccionNatural} \\
\multicolumn{4}{@{}l}{\textit{Capítulo 11: completitud}} \\
\(\bot\) en \(\{\neg,\to\}\) & \lean{evaluar_falsedad} & 3, 2, 7 & \texttt{Prop/NaturalDeduction} \\
Debilitamiento & \lean{debilitamiento} & 11 & \texttt{Prop/NaturalDeduction} \\
Reglas derivadas & \lean{introduccion_doble_negacion}, \lean{analisis_por_casos} & 13, 2 & \texttt{Prop/NaturalDeduction} \\
Corrección & \lean{correccion}, \lean{tautologia_de_derivable} & 11, 7, 1 & \texttt{Prop/NaturalDeduction} \\
Consistencia & \lean{no_derivable_falsedad} & 7, 14 & \texttt{Prop/NaturalDeduction} \\
Lema de Kalmár & \lean{kalmar} & 10, 4, 6, 7, 9, 13 & \texttt{Prop/Completeness} \\
Descarga de átomos & \lean{descarga}, \lean{congruencia_contextoLiterales} & 10, 9, 14 & \texttt{Prop/Completeness} \\
Completitud & \lean{completitud_ND}, \lean{correccion_completitud} & 14, 3 & \texttt{Prop/Completeness} \\
Validación externa & \lean{evaluar_traduccion}, \lean{traduccion_derivable_de_tautologia} & 16, 10, 6 & \texttt{Prop/Bridge…}, \texttt{Prop/CompletenessVia…} \\
\multicolumn{4}{@{}l}{\textit{Capítulo 12: Quicksort}} \\
Terminación & \lean{quicksort}, \lean{longitud_filtro_le}, \lean{longitud_filtro_lt_cons} & 12, 10 & \texttt{Sort/Quicksort} \\
Partición & \lean{particion_filtros_permutacion} & 10, 9 & \texttt{Sort/Quicksort} \\
Reglas de \lean{Pairwise} & \lean{relacion_pares_cons_iff}, \lean{relacion_pares_concatenacion_iff} & 3, 4, 10 & \texttt{Sort/Quicksort} \\
Permutación & \lean{quicksort_permutacion}, \lean{pertenencia_quicksort} & 12, 9, 14 & \texttt{Sort/Quicksort} \\
Ordenamiento & \lean{quicksort_ordenada} & 12, 5, 9, 14 & \texttt{Sort/Quicksort} \\
Corrección total & \lean{quicksort_correcto} & 4, 14 & \texttt{Sort/Quicksort} \\
\bottomrule
\end{longtable}
}

%=====================================================================
\section{Qué garantiza Lean y qué no}\label{sec:garantias}
%=====================================================================

\textbf{Lo que garantiza.} Si el archivo compila, el núcleo de Lean (un
comprobador de tipos pequeño) ha verificado que cada término de prueba tiene
el tipo del enunciado. Las tácticas no forman parte de la base de confianza:
solo construyen términos, y el núcleo los vuelve a comprobar.

\textbf{Axiomas.} El comando \texttt{\#print axioms t} lista los axiomas de los
que depende \lean{t}. La herramienta \texttt{tools/check\_axioms.py} rechaza
cualquier axioma fuera de \lean{propext}, \lean{Classical.choice} y
\lean{Quot.sound}, los tres estándar de Lean y Mathlib. También rechaza
\lean{sorryAx}, que aparecería si alguna prueba estuviera incompleta
(\lean{sorry}). Por ejemplo, la auditoría del capítulo~11 indica:
\begin{itemize}
  \item \lean{completitud_ND} y \lean{correccion} dependen de los tres axiomas:
    usan razonamiento clásico en la metalógica (\lean{by_cases},
    \lean{by_contra}).
  \item \lean{ej_identidad} y \lean{ej_eliminacion_doble_negacion} no dependen
    de ningún axioma: son construcciones directas de derivaciones. Que la
    segunda use RAA \emph{en el cálculo objeto} no requiere lógica clásica
    \emph{en Lean}.
\end{itemize}

\textbf{Lo que no garantiza.} Lean no comprueba que las definiciones
signifiquen lo que pretende el texto. Si \lean{ND} tuviera una regla de más,
la completitud sería trivial pero errónea respecto de la tesis. Por eso
importan:
\begin{itemize}
  \item la lectura cuidadosa de las definiciones (\lean{Formula}, \lean{ND},
    \lean{Deriva}, \lean{PasoBeta}, \lean{sustitucion});
  \item los teoremas de validación del esquema~17;
  \item que las hipótesis adicionales sean visibles en los enunciados
    (esquema~15) y no queden escondidas en axiomas;
  \item las diferencias declaradas en los comentarios de cabecera de cada
    archivo (sustitución ingenua con condición lateral en \texttt{Cubo.lean},
    contextos como conjuntos en \texttt{Prop}, nombres como \lean{String}).
\end{itemize}

%=====================================================================
\section{Cómo explicar una prueba en voz alta}\label{sec:explicar}
%=====================================================================

\subsection{Receta}
\begin{enumerate}
  \item \textbf{Enunciado.} Léelo en matemáticas. Señala qué está
    cuantificado, qué son hipótesis y qué es la conclusión. Si hay notación
    objeto (\lean{⟶}, \lean{~}, \lean{⊢ₚ}), dilo explícitamente.
  \item \textbf{Método.} Nombra el esquema: «por inducción sobre la
    derivación», «por casos sobre el valor de verdad», «construyendo la
    derivación línea por línea».
  \item \textbf{Generalización.} Si hay inducción, explica qué quedó
    cuantificado en la meta y por qué (esquemas~10 y~11).
  \item \textbf{Casos.} Para cada caso, di qué se sabe (hipótesis, hipótesis
    inductivas) y qué falta. Usa los estados, no las tácticas.
  \item \textbf{Paso clave.} Identifica la línea con contenido matemático
    (la aplicación de la hipótesis inductiva, la regla que descarga, la
    elección del testigo). Lo demás suele ser contabilidad.
  \item \textbf{Cierre.} Indica qué regla o lema cierra cada rama.
  \item \textbf{Confianza.} Menciona los axiomas y, si existen, las
    hipótesis explícitas.
\end{enumerate}

\subsection{Diccionario de frases}
\begin{center}
\small
\begin{tabular}{@{}ll@{}}
\toprule
\textbf{En Lean} & \textbf{Se dice} \\
\midrule
\lean{intro x hx} & «Sea \(x\) y supongamos \(hx\)\ldots» \\
\lean{have h : P := t} & «Tenemos \(P\) por \(t\)\ldots» \\
\lean{exact t} & «\ldots lo cual es exactamente \(t\)» \\
\lean{induction d with} & «Por inducción sobre la derivación \(d\)\ldots» \\
\lean{| impI … hipInd =>} & «Si la última regla es \(\to\)I, con hipótesis inductiva\ldots» \\
\lean{rcases h with a | b} & «Distinguimos dos casos\ldots» \\
\lean{by_cases h : P} & «Si \(P\)\ldots; en caso contrario\ldots» \\
\lean{by_contra h} & «Supongamos, por el absurdo, \(\neg\)(meta)\ldots» \\
\lean{obtain ⟨n, h⟩ := e} & «Sea \(n\) tal que \(h\)\ldots» \\
\lean{exact ⟨a, h⟩} & «Basta tomar \(a\); en efecto, \(h\)» \\
\lean{rw [def]} & «Por definición de \ldots» \\
\lean{calc} & «Encadenando: \(a = b = c\)» \\
\lean{ND.hyp (Set.mem_insert _ _)} & «\(\varphi\) está entre las suposiciones abiertas» \\
\lean{ND.impI d} & «Por \(\to\)I, descargando la suposición\ldots» \\
\lean{debilitamiento d h} & «Reiteramos \(d\) dentro de la caja» \\
\lean{simp [Contexto.dominio]} & «La condición lateral se comprueba directamente» \\
\lean{decreasing_by} & «La llamada es sobre una lista más corta» \\
\bottomrule
\end{tabular}
\end{center}

\subsection{Preguntas frecuentes}

\textbf{¿Por qué las pruebas son tan largas?} Porque se eligió el estilo hacia
adelante con todos los tipos intermedios escritos. Muchas se podrían acortar
con automatización (\lean{simp}, \lean{aesop}), pero entonces dejarían de
seguir la prueba en papel paso a paso. Es una decisión de legibilidad, no una
limitación.

\textbf{¿Qué hace \lean{simp}? ¿No es una caja negra?} Reescribe con lemas
marcados como simplificaciones. En la tesis se limita a desplegar definiciones
(\lean{simp only [l]}) o a cerrar condiciones decidibles sobre listas de
nombres. Además su resultado pasa por el núcleo como cualquier otro término.

\textbf{¿Por qué los contextos son conjuntos en el capítulo~11 y listas en
tipos simples?} En deducción natural importa solo qué fórmulas están
disponibles, de modo que un conjunto es lo natural y hace que el
debilitamiento sea una simple inclusión. En los sistemas de tipos la tesis
define el contexto como una lista con sujetos distintos, y el orden importa
(en \(\lambda 2\), una variable de tipo se declara antes de usarse).

\textbf{¿La lógica de Lean es clásica? ¿No afecta a lo que se demuestra?} Lean
es constructivo por defecto, y la lógica clásica entra a través del axioma
\lean{Classical.choice}. Lo que se demuestra \emph{sobre} el cálculo ND (que
toda tautología es derivable) usa razonamiento clásico en la metalógica, como
la prueba en papel. Las derivaciones \emph{dentro} de ND que usan RAA
(\lean{ND.classical}) son objetos del cálculo y no requieren axiomas en Lean.

\textbf{¿Por qué Church--Rosser es una hipótesis?} Porque la tesis lo cita sin
demostración. Declararlo como \lean{axiom} sería ocultarlo; pedirlo como
premisa \lean{hCR : Confluente sust} hace visible exactamente qué depende de
él.

\textbf{¿Cómo sé que \lean{AlfaEq} es la \(\alpha\)-equivalencia de la tesis?}
Por \lean{alfaEq_iff_alfaTesis}: coincide con la relación generada
literalmente por las reglas de la Definición~7.16.

\appendix
%=====================================================================
\section{Reproducir y explorar}
%=====================================================================

\begin{itemize}
  \item Compilar y auditar: \texttt{lake build}, \texttt{make check},
    \texttt{make audit audit-quicksort verification-output}.
  \item Ver los estados de cualquier prueba: abrir el archivo en VS Code con la
    extensión de Lean~4 y colocar el cursor tras una táctica; el panel
    \emph{Lean Infoview} muestra el estado, igual que los recuadros
    amarillos de este documento.
  \item Exportar todos los estados: \texttt{make proof-explorer} y
    \texttt{make proof-explorer-theory} regeneran
    \texttt{artifacts/explorer/proofs.json} y \texttt{theory.json}.
  \item Regenerar este documento:
    \texttt{python3 reports/curso-lean/build.py}. Si cambia el código Lean, los
    fragmentos se actualizan solos; los estados dependen de los JSON
    exportados.
\end{itemize}

\end{document}
